% AUTO-GENERATED by tools/build_crosswalk.py from tools/crosswalk_data.json -- edit the data there.
\section{Where Every Prior Category Sits in PIAtlas}
\label{app:crosswalk}
\newcommand{\xwalso}[1]{\textcolor{black!55}{\itshape partly outside PIAtlas (#1)}}

This appendix places the categories of \xwNschemes{} prior schemes in PIAtlas,
one card per scheme, first the \xwNtable{} of Table~\ref{tab:tax-compare} and then
\xwNcardonly{} more that we also checked.
The card for SaTML CTF 2024~\cite{debenedetti2024satml} follows them. That source defines no categories of its own, so its card lists only examples, and we do not count it among the \xwNschemes{} schemes. The last card places the
native labels of the ingested datasets. The title bar gives the scheme and what it organizes attacks by, and
a one-sentence summary follows. Each row pairs a category with its place in
PIAtlas, given as a class badge and the technique name, with the code in grey (see
Appendix~\ref{app:fulltax}). A category in \textit{italics} is one of the
scheme's own categories, named as the source names it. An entry in upright type
is an example, a published attack that the scheme files under one of its
categories, or description text, and never sets a filled circle in
Table~\ref{tab:tax-compare}. ``Or'' means that the placement depends on the
instance, and ``+'' means a composition. A grey \outside{goal} marks a category
defined on an axis other than construction, either one of the three that PIAtlas sets
aside (delivery, goal, and search method in Section~\ref{sec:scope}) or a
lifecycle stage, trust level, target, failure mode, or kind of model access,
which PIAtlas does not model. A category that is placed in PIAtlas but also
depends on such an axis is marked \xwalso{goal}. No category needed a seventh class. We also
read schemes that classify something other than attack text, namely agent
failure modes~\cite{microsoft2026failuremodes}, risks of agentic
applications~\cite{owasp2025agentic,owasp2025agenticthreats}, and ways of
generating attacks~\cite{shi2025gemini}. Their categories lie on other axes, so
they have no card. Rainbow Teaming's attack styles~\cite{samvelyan2024rainbow},
JailbreakRadar's obfuscation class~\cite{chu2024comprehensive}, and the
in-the-wild prompt communities of Shen et al.~\cite{shen2024dan} do include
construction categories. We have not carded them, but each falls within the
six classes. In the \emph{Multi-label} column of Table~\ref{tab:tax-compare}, \checkmark{}
means that an attack may carry several of the scheme's categories, and
\emph{facets} and \emph{stages} mean one category per facet or per stage of an
attack. The entries \emph{comb.}, \emph{comp.}, and \emph{fact.} mean that the scheme
combines its categories, builds an attack from its categories as parts, or
crosses them as factors, and a blank means one category per attack.

\paragraph*{Boundary cases}
Four kinds of category sit on the line between construction and another
axis. Rendering-level stealth (HTML comments, white or tiny text, text in
images) is delivery, because the model still receives ordinary text. Hiding
done with the characters themselves (invisible Unicode, ASCII smuggling)
reaches the tokenizer and is F3. A conditional or delayed directive
(``when the user says thanks, \ldots'') is an F1 construction that also plays a
lifecycle role. A multi-stage attack is a lifecycle pattern whose first stage,
the pointer to the rest of the payload, is fetch and follow (F1). A multi-turn
attack is placed turn by turn, and its turn structure is delivery.

\newtcolorbox{xwcard}[2]{enhanced, breakable, colback=white, colframe=black!28,
  boxrule=0.45pt, arc=1.5pt, outer arc=1.5pt,
  left=3.5pt, right=3.5pt, top=2.5pt, bottom=2.5pt,
  colbacktitle=hdrnavy, coltitle=white, toptitle=1.3pt, bottomtitle=1.3pt,
  fonttitle=\sffamily\bfseries\footnotesize,
  title={#1\unskip\nobreak\hfil\penalty50\hskip1em\hbox{}\nobreak\hfil\mbox{\mdseries\scriptsize\color{white!78!hdrnavy}#2}\parfillskip=0pt\finalhyphendemerits=0\relax},
  before skip=7pt, after skip=2pt, fontupper=\footnotesize}
\newcommand{\xwnote}[1]{{\raggedright\noindent\scriptsize\itshape\color{black!70}#1\par}\nopagebreak\vspace{2.5pt}}
\newcommand{\xwex}[1]{{\upshape #1}}
\newenvironment{xwrows}{\noindent\rowcolors{1}{black!4}{white}%
  \renewcommand{\arraystretch}{1.14}\setlength{\tabcolsep}{2.5pt}%
  \begin{tabular}{@{}>{\raggedright\arraybackslash\itshape}p{0.37\linewidth}%
  >{\raggedright\arraybackslash}p{0.60\linewidth}@{}}}{\end{tabular}}

\needspace{7\baselineskip}
\begin{xwcard}{Perez \& Ribeiro 2022~\cite{perez2022ignore}}{organizes by attack goal, plus reusable prompt blocks}
\xwnote{PromptInject divides attacks by goal, hijacking or leaking, and builds attack prompts from blocks, of which only the instruction and the delimiter blocks describe construction.}
\begin{xwrows}
goal hijacking, prompt leaking & \outside{goal} \\
instruction & (\cb{1}\,instruction nullification\,\tcode{T2.S1} \,+\, replacement command\,\tcode{T1.S2}) or proxy request\,\tcode{T7} (ignore-and-print, or spell check as proxy) \\
escape character, delimiter char, delimiter length, escape repetition & \cb{2}\,section separator\,\tcode{T1.S2} (falsely signals a new prompt section) \\
rogue string & \outside{goal}, the target text to print \\
scoring method & \outside{goal}, measures success, not attack text \\
instruction, n-shot examples, n of n-shot examples, secret instruction, private value, name of the human, name of the AI (base prompt blocks) & \outside{target}, the application prompt under attack \\
\end{xwrows}\par
\begin{xwrows}
model, temperature, top-p, presence penalty, frequency penalty, max tokens, stop sequence & \outside{target}, the model and its settings \\
\xwex{the same instruction in upper case} & \cb{1}\,capitals and symbols\,\tcode{T4.S1} \\
\xwex{instruction examples of Table C2 (print it in a single line)} & \cb{6}\,forced format\,\tcode{T3} \\
\end{xwrows}\par
\end{xwcard}

\needspace{7\baselineskip}
\begin{xwcard}{Greshake et al. 2023~\cite{greshake2023notwhat}}{organizes by injection method, threat, and affected party}
\xwnote{The paper organizes indirect injection by injection method, threat, and affected party. Construction appears only in its hidden injections and its example prompts.}
\begin{xwrows}
passive methods, active methods, user-driven injections & \outside{delivery channel} \\
hidden injections & \cb{1}\,fetch and follow\,\tcode{T6.S4} or \cb{3}\,standard encoding\,\tcode{T2.S1} or pseudocode\,\tcode{T8.S2} (multi-stage, image, encoded, or program output), \xwalso{lifecycle, delivery modality} \\
information gathering, fraud, intrusion, malware, manipulated content, availability, and their 18 sub-threats in Fig. 2 & \outside{goal, lifecycle}, persistence and worm spreading are lifecycle \\
remote control & \cb{1}\,fetch and follow\,\tcode{T6.S4} (fetches new instructions), \xwalso{goal} \\
end-users, developers, automated systems, the LLM itself & \outside{target}, affected parties \\
\xwex{example prompts in its appendix} & (\cb{2}\,role label\,\tcode{T1.S1} \,+\, status or error message\,\tcode{T5.S2} \,+\, \cb{5}\,unrestricted AI\,\tcode{T1.S1} \,+\, \cb{1}\,policy revocation\,\tcode{T2.S4} \,+\, \cb{2}\,assistant turn\,\tcode{T4.S1}) or (\cb{1}\,direct aside\,\tcode{T5.S2} \,+\, \cb{2}\,markup element\,\tcode{T3.S2} \,+\, \cb{1}\,capitals and symbols\,\tcode{T4.S1}) or (\cb{4}\,official notice\,\tcode{T1.S3} \,+\, \cb{1}\,fetch and follow\,\tcode{T6.S4} \,+\, conditional trigger\,\tcode{T6.S1} \,+\, \cb{6}\,fixed text\,\tcode{T3.S2}) or (\cb{2}\,chat-template token\,\tcode{T2.S1} \,+\, \cb{5}\,story or script\,\tcode{T3.S1} \,+\, unrestricted AI\,\tcode{T1.S1}) or (\cb{2}\,chat-template token\,\tcode{T2.S1} \,+\, \cb{5}\,unrestricted AI\,\tcode{T1.S1} \,+\, \cb{3}\,variable concatenation\,\tcode{T4.S1}) or (\cb{2}\,chat-template token\,\tcode{T2.S1} \,+\, assistant turn\,\tcode{T4.S1} \,+\, status or error message\,\tcode{T5.S2} \,+\, \cb{4}\,emergency\,\tcode{T3.S2} \,+\, \cb{1}\,fetch and follow\,\tcode{T6.S4}) or (\cb{2}\,role label\,\tcode{T1.S1} \,+\, \cb{3}\,standard encoding\,\tcode{T2.S1}) (one alternative per prompt family) \\
\end{xwrows}\par
\end{xwcard}

\needspace{7\baselineskip}
\begin{xwcard}{Liu et al. 2023~\cite{liu2023jailbreaking}}{organizes by prompt pattern, grouped into three types}
\xwnote{The study sorts jailbreak prompts into ten patterns grouped in three types. Most patterns describe how the prompt is built, such as personas, pretexts, or continuation.}
\begin{xwrows}
pretending & \cb{5}\,scenario framing (changes context, keeps intention) \\
attention shifting & \cb{6}\,seeded continuation\,\tcode{T1.S2} or \cb{4}\,argument and obligation\,\tcode{T7} or \cb{3}\,pseudocode\,\tcode{T8.S2} or \cb{1}\,transformation object\,\tcode{T3.S3} or \cb{3}\,foreign-language payload\,\tcode{T3.S1} (depends on the pattern) \\
privilege escalation & \cb{5}\,persona\,\tcode{T1} or \cb{4}\,authority claim\,\tcode{T1} (claimed special mode or authority) \\
character role play & \cb{5}\,persona\,\tcode{T1} \\
assumed responsibility & \cb{4}\,responsibility shift\,\tcode{T7.S5} or duty to help\,\tcode{T7.S4} (responsibility shift or duty to help) \\
research experiment & \cb{4}\,research purpose\,\tcode{T4.S1} or \cb{5}\,story or script\,\tcode{T3.S1} (stated study, or a fictional lab scene) \\
\end{xwrows}\par
\begin{xwrows}
text continuation & \cb{6}\,seeded continuation\,\tcode{T1.S2} or (seeded continuation\,\tcode{T1.S2} \,+\, \cb{5}\,story or script\,\tcode{T3.S1}) (often a story to continue) \\
logical reasoning & \cb{4}\,argument and obligation\,\tcode{T7} (no example given in the paper) \\
program execution & \cb{3}\,pseudocode\,\tcode{T8.S2} (the prompt is written as a program) \\
translation & \cb{1}\,transformation object\,\tcode{T3.S3} or (transformation object\,\tcode{T3.S3} \,+\, \cb{3}\,foreign-language payload\,\tcode{T3.S1}) (a translation task, with or without a foreign payload) \\
superior model & \cb{5}\,unrestricted AI\,\tcode{T1.S1} or \cb{4}\,insider identity\,\tcode{T1.S1} (closest match, as the definition is vague) \\
sudo mode & \cb{5}\,special mode\,\tcode{T1.S2} \\
\end{xwrows}\par
\begin{xwrows}
simulate jailbreaking & \cb{5}\,unrestricted AI\,\tcode{T1.S1} or special mode\,\tcode{T1.S2} (jailbroken AI or developer mode) \\
illegal activities, harmful content, fraudulent or deceptive activities, adult content, political campaigning or lobbying, violating privacy, unlawful practices, high-risk government decision-making & \outside{goal}, prohibited scenarios used for testing \\
\xwex{example prompts in Secs. II-B and IV-A} & (\cb{5}\,story or script\,\tcode{T3.S1} \,+\, \cb{6}\,seeded continuation\,\tcode{T1.S2}) or (\cb{3}\,pseudocode\,\tcode{T8.S2} \,+\, \cb{1}\,sample request\,\tcode{T7.S1}) or (\cb{5}\,special mode\,\tcode{T1.S2} \,+\, \cb{1}\,policy revocation\,\tcode{T2.S4}) \\
\end{xwrows}\par
\end{xwcard}

\needspace{7\baselineskip}
\begin{xwcard}{Wei et al. 2023~\cite{wei2023jailbroken}}{organizes by failure mode of safety training}
\xwnote{The paper explains jailbreaks by two failure modes of safety training and builds most of its 30 tested jailbreak methods from them, mostly output coercion and obfuscation.}
\begin{xwrows}
competing objectives, mismatched generalization & \outside{failure mode}, why attacks work, not how built \\
prefix\_injection, prefix\_injection\_hello & \cb{6}\,required opening\,\tcode{T1.S1} (the second is a benign-prefix ablation) \\
refusal\_suppression, refusal\_suppression\_inv & \cb{6}\,apology and disclaimer ban\,\tcode{T2.S1} or style mandate\,\tcode{T3.S5} (the inverted ablation only dictates style) \\
style\_injection\_short, style\_injection\_json & \cb{6}\,style mandate\,\tcode{T3.S5} or schema only\,\tcode{T3.S1} \\
base64, base64\_input\_only, base64\_raw & (\cb{3}\,standard encoding\,\tcode{T2.S1} \,+\, \cb{6}\,altered answer\,\tcode{T3.S3}) or \cb{3}\,standard encoding\,\tcode{T2.S1} (base64 also asks for an encoded answer) \\
base64\_output\_only & \cb{6}\,altered answer\,\tcode{T3.S3} (plain question, encoded answer) \\
\end{xwrows}\par
\begin{xwrows}
disemvowel, leetspeak, rot13 & \cb{3}\,character perturbation\,\tcode{T1} or leetspeak\,\tcode{T1.S1} or classical cipher\,\tcode{T2.S2} \\
distractors, poems & \cb{3}\,distractors\,\tcode{T5.S5} or (distractors\,\tcode{T5.S5} \,+\, \cb{6}\,style mandate\,\tcode{T3.S5}) (poems asks for unrelated verse) \\
distractors\_negated & \cb{3}\,distractors\,\tcode{T5.S5} \,+\, \cb{1}\,instruction nullification\,\tcode{T2.S1} (the distractor is then revoked) \\
few\_shot\_json & \cb{1}\,worked examples\,\tcode{T5.S4} \,+\, \cb{6}\,pattern continuation\,\tcode{T1.S3} \,+\, schema only\,\tcode{T3.S1} \,+\, \cb{4}\,harm-prevention purpose\,\tcode{T4.S3} (JSON pairs with a classifier-dataset pretext) \\
wikipedia, wikipedia\_with\_title & (\cb{1}\,composition request\,\tcode{T7.S2} \,+\, \cb{6}\,style mandate\,\tcode{T3.S5}) or (\cb{1}\,composition request\,\tcode{T7.S2} \,+\, \cb{6}\,style mandate\,\tcode{T3.S5} \,+\, required opening\,\tcode{T1.S1}) (the second adds a title opening) \\
combination\_1, combination\_2, combination\_3 & (\cb{6}\,required opening\,\tcode{T1.S1} \,+\, apology and disclaimer ban\,\tcode{T2.S1} \,+\, \cb{3}\,standard encoding\,\tcode{T2.S1} \,+\, \cb{6}\,altered answer\,\tcode{T3.S3}) or (required opening\,\tcode{T1.S1} \,+\, apology and disclaimer ban\,\tcode{T2.S1} \,+\, \cb{3}\,standard encoding\,\tcode{T2.S1} \,+\, \cb{6}\,altered answer\,\tcode{T3.S3} \,+\, style mandate\,\tcode{T3.S5}) or (required opening\,\tcode{T1.S1} \,+\, apology and disclaimer ban\,\tcode{T2.S1} \,+\, \cb{3}\,standard encoding\,\tcode{T2.S1} \,+\, \cb{6}\,altered answer\,\tcode{T3.S3} \,+\, style mandate\,\tcode{T3.S5} \,+\, \cb{1}\,composition request\,\tcode{T7.S2}) (each adds to the previous one) \\
\end{xwrows}\par
\begin{xwrows}
auto\_payload\_splitting & \cb{3}\,variable concatenation\,\tcode{T4.S1} (GPT-4 picks the terms to split), \xwalso{search method} \\
auto\_obfuscation & \cb{3}\,homoglyphs\,\tcode{T1.S3} or (leetspeak\,\tcode{T1.S1} \,+\, foreign-language payload\,\tcode{T3.S1}) or standard encoding\,\tcode{T2.S1} (the model invents the obfuscation), \xwalso{search method} \\
AIM, dev\_mode\_v2, evil\_confidant & (\cb{5}\,unrestricted AI\,\tcode{T1.S1} \,+\, \cb{6}\,required opening\,\tcode{T1.S1}) or (\cb{5}\,special mode\,\tcode{T1.S2} \,+\, \cb{1}\,policy revocation\,\tcode{T2.S4} \,+\, \cb{6}\,required opening\,\tcode{T1.S1}) or (\cb{5}\,villain persona\,\tcode{T1.S3} \,+\, \cb{1}\,policy revocation\,\tcode{T2.S4} \,+\, \cb{6}\,required opening\,\tcode{T1.S1}) (role play plus forced response prefixes) \\
dev\_mode\_with\_rant & \cb{5}\,special mode\,\tcode{T1.S2} \,+\, \cb{1}\,policy revocation\,\tcode{T2.S4} \,+\, \cb{6}\,staged denial\,\tcode{T1.S5} (a rant first, then compliance) \\
evil\_system\_prompt & \cb{5}\,villain persona\,\tcode{T1.S3} \,+\, \cb{1}\,policy revocation\,\tcode{T2.S4} \,+\, \cb{6}\,required opening\,\tcode{T1.S1} (Evil Confidant placed in the system prompt), \xwalso{delivery channel} \\
adaptive attack & \outside{search method}, succeeds if any attack succeeds \\
\end{xwrows}\par
\begin{xwrows}
\xwex{other examples named in Sec. 3 (Morse code, Pig Latin, synonyms, translation, JSON output, DAN, a rant before complying)} & \cb{3}\,classical cipher\,\tcode{T2.S2} or synonym swap\,\tcode{T7.S1} or foreign-language payload\,\tcode{T3.S1} or \cb{6}\,schema only\,\tcode{T3.S1} or (\cb{5}\,unrestricted AI\,\tcode{T1.S1} \,+\, \cb{6}\,required opening\,\tcode{T1.S1}) or staged denial\,\tcode{T1.S5} \\
\end{xwrows}\par
\end{xwcard}

\needspace{7\baselineskip}
\begin{xwcard}{HackAPrompt 2023~\cite{schulhoff2023hackaprompt}}{organizes by technique, linked by kind-of and part-of relations}
\xwnote{The ontology links 29 prompt-hacking techniques by kind-of and part-of relations. Most describe how the attack text is built, while intents and delivery are listed separately.}
\begin{xwrows}
simple instruction, compound instruction & \cb{1}\,plain command\,\tcode{T1} \\
context ignoring & \cb{1}\,instruction nullification\,\tcode{T2.S1} \\
special case & \cb{1}\,exception clause\,\tcode{T2.S5} or importance flag\,\tcode{T4.S2} (defined by a special-case statement) \\
instruction repetition & \cb{1}\,repetition\,\tcode{T4.S3} \\
distractor instructions, negated distractor instructions & \cb{3}\,distractors\,\tcode{T5.S5} \,+\, \cb{1}\,instruction nullification\,\tcode{T2.S1} (extra tasks, then revoked) \\
refusal suppression & \cb{6}\,apology and disclaimer ban\,\tcode{T2.S1} \\
\end{xwrows}\par
\begin{xwrows}
style injection, prefix injection, text completion as instruction & \cb{6}\,style mandate\,\tcode{T3.S5} or required opening\,\tcode{T1.S1} or seeded continuation\,\tcode{T1.S2} (prefix injection is in App. D only) \\
few shot & \cb{1}\,worked examples\,\tcode{T5.S4} (input-output pairs without role labels) \\
defined dictionary & \cb{1}\,meaning remap\,\tcode{T6.S5} \,+\, \cb{3}\,pseudocode\,\tcode{T8.S2} \,+\, \cb{1}\,worked examples\,\tcode{T5.S4} (a Python dict remaps the closing instruction) \\
cognitive hacking & \cb{5}\,persona\,\tcode{T1} (role prompting) \\
virtualization & \cb{5}\,persona\,\tcode{T1} or (persona\,\tcode{T1} \,+\, \cb{4}\,foot-in-the-door\,\tcode{T5.S1}) (gradual shift over several prompts), \xwalso{delivery turn structure} \\
task deflection & \cb{1}\,proxy request\,\tcode{T7} (the task is masked as another) \\
\end{xwrows}\par
\begin{xwrows}
fill in the blank, payload splitting, variable assignment & \cb{3}\,fragmentation\,\tcode{T4} or variable concatenation\,\tcode{T4.S1} \\
code injection & \cb{3}\,pseudocode\,\tcode{T8.S2} (Fig. 5 node without a text definition) \\
context overflow & \cb{3}\,filler padding\,\tcode{T5.S1} (padding that limits the output length) \\
context switching & \cb{1}\,task extension\,\tcode{T3.S1} or \cb{2}\,section separator\,\tcode{T1.S2} or assistant turn\,\tcode{T4.S1} or \cb{5}\,scenario framing (hypernym of continuation and termination) \\
context continuation & \cb{1}\,task extension\,\tcode{T3.S1} \\
context termination, separators & \cb{2}\,assistant turn\,\tcode{T4.S1} or section separator\,\tcode{T1.S2} (separators are parts, not attacks) \\
\end{xwrows}\par
\begin{xwrows}
obfuscation, syntactic transformation, typos & \cb{3}\,standard encoding\,\tcode{T2.S1} or classical cipher\,\tcode{T2.S2} or leetspeak\,\tcode{T1.S1} or case and typo noise\,\tcode{T1.S6} (Base64, ROT13 or Pig Latin, leet, typos) \\
translation & \cb{3}\,foreign-language payload\,\tcode{T3.S1} \\
anomalous token & \cb{3}\,non-semantic sequence\,\tcode{T6} \\
recursive prompt hacking & \outside{target}, a second, checking model is the target \\
explicit instructions, implicit instructions & \cb{1}\,plain command\,\tcode{T1} or indirect phrasing\,\tcode{T5} (plain command versus implicit instruction) \\
direct prompt injection, indirect prompt injection & \outside{delivery channel} \\
\end{xwrows}\par
\begin{xwrows}
prompt leaking, training data reconstruction, malicious action generation, harmful information generation, token wasting, denial of service, token theft & \outside{goal}, intents \\
competing objectives, mismatched generalization & \outside{failure mode} \\
\xwex{example prompts in Sec. 5 and Appendix D} & (\cb{2}\,assistant turn\,\tcode{T4.S1} \,+\, \cb{5}\,game objective\,\tcode{T4.S1}) or (terminal or interpreter\,\tcode{T2.S1} \,+\, \cb{3}\,variable concatenation\,\tcode{T4.S1} \,+\, pseudocode\,\tcode{T8.S2}) or \cb{1}\,composition request\,\tcode{T7.S2} (termination, recursive, and deflection examples) \\
target phrase generation & \outside{goal} \\
\end{xwrows}\par
\end{xwcard}

\needspace{7\baselineskip}
\begin{xwcard}{Kang et al. 2024~\cite{kang2024exploiting}}{organizes by classical program attacks adapted to prompts}
\xwnote{The paper adapts three classical program attacks to prompts (obfuscation, code injection with payload splitting, and virtualization), and all three describe how the prompt is built.}
\begin{xwrows}
obfuscation & \cb{3}\,case and typo noise\,\tcode{T1.S6} or synonym swap\,\tcode{T7.S1} (typos or synonyms, with encoding only conjectured) \\
code injection / payload splitting & \cb{3}\,variable concatenation\,\tcode{T4.S1} \,+\, pseudocode\,\tcode{T8.S2} (strings as variables, joined in code) \\
virtualization & \cb{5}\,story or script\,\tcode{T3.S1} or (story or script\,\tcode{T3.S1} \,+\, \cb{4}\,foot-in-the-door\,\tcode{T5.S1}) (setup prompts may secure commitment), \xwalso{delivery turn structure} \\
hate speech, conspiracy theory promotion, phishing attacks, scams, product astroturfing & \outside{goal}, the malicious uses it tests \\
\xwex{framing inside its prompt templates (SmartGPT, writing as a character)} & \cb{5}\,unrestricted AI\,\tcode{T1.S1} or villain persona\,\tcode{T1.S3} \\
\xwex{capability demonstrations in Sec. 2 (string concatenation, variable assignment, branching)} & \cb{3}\,variable concatenation\,\tcode{T4.S1} or \cb{1}\,conditional trigger\,\tcode{T6.S1} (capabilities, not attack classes) \\
\end{xwrows}\par
\end{xwcard}

\needspace{7\baselineskip}
\begin{xwcard}{Rao et al. 2024~\cite{rao2024tricking}}{organizes by linguistic level of the technique, plus intent}
\xwnote{It classifies each jailbreak input by the linguistic level at which it is built, from spelling to pragmatics, and separately by the attacker's intent.}
\begin{xwrows}
orthographic techniques (ORTH), syntactical transformation (SYN) & \cb{3}\,character perturbation\,\tcode{T1} or encoding\,\tcode{T2} (leetspeak, Base64, transliteration) \\
lexical techniques & \cb{3}\,non-semantic sequence\,\tcode{T6} or distractors\,\tcode{T5.S5} (e.g., distractors (Sec. 4.3)) \\
morpho-syntactic techniques, text completion as instruction (TCINS) & \cb{6}\,seeded continuation\,\tcode{T1.S2} \\
semantic techniques & (\cb{1}\,instruction nullification\,\tcode{T2.S1} \,+\, replacement command\,\tcode{T1.S2}) or worked examples\,\tcode{T5.S4} or \cb{6}\,refusal suppression\,\tcode{T2} (e.g., direct instruction, few-shot, refusal suppression) \\
direct instruction (INSTR) & \cb{1}\,instruction nullification\,\tcode{T2.S1} \,+\, replacement command\,\tcode{T1.S2} \\
few-shot hacking (FSH) & \cb{1}\,worked examples\,\tcode{T5.S4} (unlabeled input-output examples) \\
\end{xwrows}\par
\begin{xwrows}
pragmatic techniques & \cb{1}\,repetition\,\tcode{T4.S3} or proxy request\,\tcode{T7} or \cb{4}\,social engineering or \cb{5}\,scenario framing or \cb{6}\,style mandate\,\tcode{T3.S5} (persuasion, implicature, role-play, and sub-types) \\
instruction repetition (IR) & \cb{1}\,repetition\,\tcode{T4.S3} or (repetition\,\tcode{T4.S3} \,+\, \cb{4}\,sympathy\,\tcode{T2.S1}) (pleading phrases may be added) \\
indirect task deflection (ITD) & \cb{1}\,proxy request\,\tcode{T7} (its example asks for code (S2)) \\
cognitive hacking (COG) & \cb{5}\,persona\,\tcode{T1} or consequence-free sandbox\,\tcode{T2.S3} or fiction and hypotheticals\,\tcode{T3} (DAN-style personas or safe-space scenarios) \\
information leakage, misaligned content generation, performance degradation & \outside{goal} \\
prompt leaking, goal hijacking, denial of service & \outside{goal}, named sub-types of the intents \\
\end{xwrows}\par
\begin{xwrows}
user, man-in-the-middle (MITM) & \outside{delivery}, jailbreak modes \\
\xwex{its combined example (ignore, murderer persona, completion)} & \cb{1}\,instruction nullification\,\tcode{T2.S1} \,+\, \cb{5}\,villain persona\,\tcode{T1.S3} \,+\, \cb{6}\,seeded continuation\,\tcode{T1.S2} \\
\xwex{role-labeled exchanges and an assistant-voiced refusal in its jailbreak tables} & \cb{2}\,fabricated dialogue\,\tcode{T4} or assistant turn\,\tcode{T4.S1} \\
\end{xwrows}\par
\end{xwcard}

\needspace{7\baselineskip}
\begin{xwcard}{Tensor Trust 2024~\cite{toyer2024tensortrust}}{organizes by goal, plus attack strategies from topic modeling}
\xwnote{The paper splits attacks by goal, hijacking or extraction, and names sixteen attack strategies found by topic modeling, nearly all of which describe how the attack is built.}
\begin{xwrows}
prompt hijacking, prompt extraction & \outside{goal} \\
ask directly & \cb{1}\,plain command\,\tcode{T1} \\
update instructions & \cb{1}\,instruction nullification\,\tcode{T2.S1} or added command\,\tcode{T1.S1} (ignore old rules or add new ones) \\
end ambiguity & (\cb{2}\,section separator\,\tcode{T1.S2} \,+\, \cb{1}\,override\,\tcode{T2}) or override\,\tcode{T2} (marker line or prose claim) \\
few-shot & \cb{2}\,many-shot demonstrations\,\tcode{T4.S2} (role-labeled example exchanges) \\
roleplay & \cb{5}\,persona\,\tcode{T1} \\
\end{xwrows}\par
\begin{xwrows}
sentient door & \cb{5}\,persona\,\tcode{T1} \,+\, \cb{1}\,conditional trigger\,\tcode{T6.S1} (door persona with a new trigger phrase) \\
creative writing & \cb{5}\,story or script\,\tcode{T3.S1} (poem, rap, or story holding the password) \\
response prefix & \cb{6}\,staged denial\,\tcode{T1.S5} or required opening\,\tcode{T1.S1} (usually opens with its standard rejection) \\
write code & \cb{6}\,altered answer\,\tcode{T3.S3} or \cb{1}\,composition request\,\tcode{T7.S2} (e.g., the instructions restated as pseudocode) \\
execute code & \cb{3}\,pseudocode\,\tcode{T8.S2} \,+\, \cb{1}\,verbatim execution\,\tcode{T6.S2} (a print statement to run) \\
binary & \cb{3}\,standard encoding\,\tcode{T2.S1} \\
\end{xwrows}\par
\begin{xwrows}
no spaces & \cb{3}\,spacing changes\,\tcode{T1.S2} \\
artisanlib & \cb{3}\,non-semantic sequence\,\tcode{T6} (rare glitch token) \\
repeated characters & \cb{3}\,non-semantic sequence\,\tcode{T6} or filler padding\,\tcode{T5.S1} or \cb{2}\,section separator\,\tcode{T1.S2} (noise block or separator line) \\
check understanding & \cb{1}\,routine step\,\tcode{T3.S2} or proxy request\,\tcode{T7} (repeat instructions to confirm comprehension) \\
access code placeholder & \cb{4}\,false premise\,\tcode{T7.S8} or \cb{3}\,allusion\,\tcode{T7.S3} (closest match, which names the code, not its value) \\
\xwex{example attacks in Table A7 and Appendix J} & (\cb{5}\,story or script\,\tcode{T3.S1} \,+\, \cb{2}\,section separator\,\tcode{T1.S2} \,+\, \cb{4}\,threat\,\tcode{T3.S4}) or (\cb{2}\,many-shot demonstrations\,\tcode{T4.S2} \,+\, \cb{4}\,threat\,\tcode{T3.S4}) or (\cb{5}\,persona\,\tcode{T1} \,+\, \cb{1}\,conditional trigger\,\tcode{T6.S1} \,+\, \cb{2}\,user turn\,\tcode{T4.S4}) (roleplay, few-shot, and sentient-door examples) \\
\end{xwrows}\par
\begin{xwrows}
\xwex{confusing the model about the preceding prompt} & \cb{1}\,override\,\tcode{T2} or \cb{3}\,filler padding\,\tcode{T5.S1} or \cb{2}\,chat-template token\,\tcode{T2.S1} or \cb{3}\,non-semantic sequence\,\tcode{T6} (narrative in Sec. 4.2) \\
\end{xwrows}\par
\end{xwcard}

\needspace{7\baselineskip}
\begin{xwcard}{Rossi et al. 2024~\cite{rossi2024early}}{organizes by delivery (direct or indirect), then class}
\xwnote{The paper splits prompt injections into direct and indirect ones. Its six direct classes describe how the prompt is built, and its four indirect classes mostly describe delivery.}
\begin{xwrows}
direct prompt injections, indirect prompt injections & \outside{delivery channel} \\
double character & \cb{5}\,split persona\,\tcode{T1.S4} or (split persona\,\tcode{T1.S4} \,+\, \cb{6}\,forced format\,\tcode{T3}) (two responses, one unrestricted) \\
virtualization & \cb{5}\,special mode\,\tcode{T1.S2} or unrestricted AI\,\tcode{T1.S1} or simulation\,\tcode{T2} or fiction and hypotheticals\,\tcode{T3} (developer mode, free AI, VM, or scenario) \\
obfuscation & \cb{3}\,standard encoding\,\tcode{T2.S1} or synonym swap\,\tcode{T7.S1} or case and typo noise\,\tcode{T1.S6} (Base64, synonyms, or typos) \\
payload splitting & \cb{3}\,fragmentation\,\tcode{T4} (parts spread over several prompts), \xwalso{delivery turn structure} \\
adversarial suffix & \cb{3}\,non-semantic sequence\,\tcode{T6} (computed by an optimizer), \xwalso{search method} \\
\end{xwrows}\par
\begin{xwrows}
instruction manipulation & \cb{1}\,instruction nullification\,\tcode{T2.S1} or plain command\,\tcode{T1} (ignore the instructions, or ask for them), \xwalso{goal} \\
active injections, passive injections, user-driven injections & \outside{delivery channel} \\
virtual prompt injection & \outside{model access}, poisons instruction-tuning data \\
\xwex{examples in Appendix A (Table 4)} & (\cb{5}\,villain persona\,\tcode{T1.S3} \,+\, story or script\,\tcode{T3.S1} \,+\, \cb{4}\,sympathy\,\tcode{T2.S1}) or (\cb{5}\,job persona\,\tcode{T1.S6} \,+\, \cb{1}\,transformation object\,\tcode{T3.S3}) or \cb{5}\,alternate world\,\tcode{T3.S3} or (story or script\,\tcode{T3.S1} \,+\, \cb{6}\,forced format\,\tcode{T3}) (grandma, copy writer, opposite mode, Tom and Jerry) \\
\end{xwrows}\par
\end{xwcard}

\needspace{7\baselineskip}
\begin{xwcard}{HouYi 2023~\cite{liu2024houyi}}{organizes by components of the injected prompt}
\xwnote{HouYi builds an injection from framework, separator and disruptor components. Construction appears in its separator and disruptor strategies and in its categories of earlier attacks.}
\begin{xwrows}
direct injection & \cb{1}\,added command\,\tcode{T1.S1} \\
escape characters, syntax-based strategy & \cb{2}\,section separator\,\tcode{T1.S2} \\
context ignoring, specific ignoring & \cb{1}\,instruction nullification\,\tcode{T2.S1} \\
reasoning summary & \cb{1}\,task extension\,\tcode{T3.S1} (a follow-up on the task just done) \\
additional task & \cb{1}\,added command\,\tcode{T1.S1} or task extension\,\tcode{T3.S1} \\
language switching & \cb{3}\,mid-sentence switch\,\tcode{T3.S2} \\
\end{xwrows}\par
\begin{xwrows}
format & \cb{6}\,forced format\,\tcode{T3} (to match the application's output) \\
short length & \cb{6}\,partial answer\,\tcode{T3.S4} \\
framework component, reproducible answer, short answer & \outside{target}, benign request matching the application \\
disruptor component, exploit scenarios (prompt leaking, code generation, content manipulation, spam generation, information gathering) & \outside{goal} \\
application context inference, prompt refinement with dynamic feedback & \outside{search method} \\
\xwex{example injection against DecisionAI} & \cb{2}\,section separator\,\tcode{T1.S2} \,+\, \cb{1}\,instruction nullification\,\tcode{T2.S1} \,+\, \cb{3}\,mid-sentence switch\,\tcode{T3.S2} \,+\, \cb{6}\,forced format\,\tcode{T3} \\
\end{xwrows}\par
\end{xwcard}

\needspace{7\baselineskip}
\begin{xwcard}{Open-Prompt-Injection 2024~\cite{liu2024openpi}}{organizes by how the compromised data is composed}
\xwnote{Open-Prompt-Injection formalizes an attack as crafted compromised data, and each of its five attack types is a directive, a structural trick, or both.}
\begin{xwrows}
naive attack & \cb{1}\,added command\,\tcode{T1.S1} \\
escape characters & \cb{2}\,section separator\,\tcode{T1.S2} \\
context ignoring & \cb{1}\,instruction nullification\,\tcode{T2.S1} \\
fake completion & \cb{2}\,assistant turn\,\tcode{T4.S1} \\
combined attack & \cb{2}\,section separator\,\tcode{T1.S2} \,+\, assistant turn\,\tcode{T4.S1} \,+\, \cb{1}\,instruction nullification\,\tcode{T2.S1} \\
\end{xwrows}\par
\end{xwcard}

\needspace{7\baselineskip}
\begin{xwcard}{Zeng et al. 2024~\cite{zeng2024johnny}}{organizes by persuasion technique from social science}
\xwnote{Zeng et al. derive forty persuasion techniques from social-science research. Each supplies a reason to comply (F4), and the attack prompts built from them also reword the request and sometimes add a persona or fiction.}
\begin{xwrows}
expert endorsement, authority endorsement & \cb{4}\,endorsement\,\tcode{T1.S7} \\
misrepresentation & \cb{4}\,expert credentials\,\tcode{T1.S2} or pretext\,\tcode{T4} (a false self or a false issue) \\
complimenting, encouragement, affirmation, positive emotional appeal, negative emotional appeal, storytelling, priming, compensation, discouragement & \cb{4}\,emotional appeal\,\tcode{T2} \\
supply scarcity, time pressure, threats, rumors, exploiting weakness & \cb{4}\,urgency and threat\,\tcode{T3} \\
foot-in-the-door, public commitment, reciprocity & \cb{4}\,commitment\,\tcode{T5} (its reciprocity means mirroring the other party) \\
social proof, injunctive norm, non-expert testimonial, alliance building, shared values, relationship leverage, loyalty appeals, social punishment & \cb{4}\,social proof\,\tcode{T6} \\
\end{xwrows}\par
\begin{xwrows}
evidence-based persuasion, logical appeal, door-in-the-face, favor, negotiation, anchoring, framing, confirmation bias, reflective thinking, false promises, false information, creating dependency & \cb{4}\,argument and obligation\,\tcode{T7} \\
\xwex{example persuasive prompt in its Figure 2} & \cb{4}\,emotional appeal\,\tcode{T2} \,+\, research purpose\,\tcode{T4.S1} \,+\, harm-prevention purpose\,\tcode{T4.S3} \\
\xwex{its persuasive paraphrases of the plain query} & \cb{3}\,rewording\,\tcode{T7} \,+\, \cb{4}\,social engineering \\
\xwex{its attack prompts that ask for the output as an example for another purpose} & \cb{1}\,sample request\,\tcode{T7.S1} \\
\xwex{its attack prompts that assign a persona, a virtual scenario, or fiction} & \cb{5}\,persona\,\tcode{T1} or simulation\,\tcode{T2} or fiction and hypotheticals\,\tcode{T3} \\
\end{xwrows}\par
\end{xwcard}

\needspace{7\baselineskip}
\begin{xwcard}{Instruction hierarchy 2024~\cite{wallace2024hierarchy}}{organizes by privilege of the message source, plus attack goal}
\xwnote{The instruction hierarchy ranks messages by source privilege and sorts attacks by delivery and goal. Construction appears only in its example prompts.}
\begin{xwrows}
system messages, user messages, model outputs, tool outputs & \outside{trust}, from highest to lowest privilege \\
direct prompt injections, indirect prompt injections & \outside{delivery} \\
jailbreaks, system message extraction & \outside{goal} \\
aligned instructions, misaligned instructions & \outside{trust, goal}, relative to higher-privileged instructions \\
open-domain tasks, closed-domain tasks & \outside{target}, the kind of application attacked \\
\xwex{example prompts that say IGNORE PREVIOUS INSTRUCTIONS} & \cb{1}\,instruction nullification\,\tcode{T2.S1} \,+\, capitals and symbols\,\tcode{T4.S1} \\
\end{xwrows}\par
\begin{xwrows}
\xwex{example headers that announce new instructions or a system note} & \cb{2}\,section separator\,\tcode{T1.S2} or role label\,\tcode{T1.S1} \\
\xwex{example role and mode changes (DAN, a greeting bot, safe mode off)} & \cb{5}\,unrestricted AI\,\tcode{T1.S1} or job persona\,\tcode{T1.S6} or special mode\,\tcode{T1.S2} \\
\xwex{example pleas and claimed system upgrades} & \cb{4}\,sympathy\,\tcode{T2.S1} or official notice\,\tcode{T1.S3} \\
\xwex{example rule that fires on every hobby request} & \cb{1}\,conditional trigger\,\tcode{T6.S1} \\
\xwex{example demands to output a given text and nothing else} & \cb{6}\,fixed text\,\tcode{T3.S2} \\
\end{xwrows}\par
\end{xwcard}

\needspace{7\baselineskip}
\begin{xwcard}{CyberSecEval 2024~\cite{meta2024cyberseceval}}{organizes by injection method (15), type, and risk}
\xwnote{Labels each prompt-injection test case with one of fifteen injection methods, which are constructions, and separately with its delivery type and risk category.}
\begin{xwrows}
ignore previous instructions & \cb{1}\,instruction nullification\,\tcode{T2.S1} \\
indirect references & \cb{1}\,proxy request\,\tcode{T7} or \cb{3}\,synonym swap\,\tcode{T7.S1} \\
token smuggling & \cb{3}\,encoding\,\tcode{T2} \\
system mode & \cb{4}\,insider identity\,\tcode{T1.S1} or (insider identity\,\tcode{T1.S1} \,+\, \cb{5}\,special mode\,\tcode{T1.S2}) or \cb{2}\,role label\,\tcode{T1.S1} or markup element\,\tcode{T3.S2} (a claim in prose, or system tags in some test cases) \\
different user input language & \cb{3}\,foreign-language payload\,\tcode{T3.S1} \\
information overload & \cb{3}\,overload\,\tcode{T5} \\
\end{xwrows}\par
\begin{xwrows}
few-shot attack, many-shot attack & \cb{1}\,worked examples\,\tcode{T5.S4} or \cb{2}\,many-shot demonstrations\,\tcode{T4.S2} (Q and A pairs in the test cases, many-shot as in Anil et al.) \\
repeated-token attack & \cb{3}\,non-semantic sequence\,\tcode{T6} or filler padding\,\tcode{T5.S1} or \cb{1}\,repetition\,\tcode{T4.S3} \\
output formatting manipulation & \cb{6}\,altered answer\,\tcode{T3.S3} or partial answer\,\tcode{T3.S4} \\
hypothetical scenario & \cb{5}\,hypothetical\,\tcode{T3.S2} \\
payload splitting & \cb{3}\,fragmentation\,\tcode{T4} \\
persuasion & \cb{4}\,social engineering \\
\end{xwrows}\par
\begin{xwrows}
virtualization & \cb{5}\,fiction and hypotheticals\,\tcode{T3} \\
mixed techniques & a composition of the methods above \\
direct vs.\ indirect injection, logic-violating vs.\ security-violating risk & \outside{delivery, goal} \\
\end{xwrows}\par
\end{xwcard}

\needspace{7\baselineskip}
\begin{xwcard}{Attack Atlas 2024~\cite{rawat2024attackatlas}}{organizes by prompt syntax, form, and semantics}
\xwnote{It sorts single-turn prompt attacks only by the prompt's form and meaning, leaving out generation method and harm, so nearly all its nodes are construction.}
\begin{xwrows}
direct instructions & \cb{1}\,plain command\,\tcode{T1} \\
encoded interactions & \cb{3}\,obfuscation or \cb{6}\,forced format\,\tcode{T3} (encodings, styles, typography, formatting) \\
payload-splitting & \cb{3}\,fragmentation\,\tcode{T4} \\
output encoding, enforce style & \cb{6}\,forced format\,\tcode{T3} (enforce style is a style mandate (S5)) \\
surrogate modality & \cb{3}\,surrogate syntax\,\tcode{T8} or \cb{6}\,schema only\,\tcode{T3.S1} (request as JSON, CSV, or code) \\
typos/misspellings & \cb{3}\,case and typo noise\,\tcode{T1.S6} \\
\end{xwrows}\par
\begin{xwrows}
nesting & \cb{1}\,task embedding\,\tcode{T3} (request folded into another task) \\
obfuscation & \cb{3}\,obfuscation \\
pseudonym & \cb{3}\,allusion\,\tcode{T7.S3} or variable concatenation\,\tcode{T4.S1} (coded words or variables) \\
word games/puzzles & \cb{3}\,cognitive overload\,\tcode{T5.S2} (the puzzle's answer carries the goal) \\
token smuggling & \cb{3}\,encoding\,\tcode{T2} (ASCII, Base64, or Morse) \\
low-resource, uncommon dialect & \cb{3}\,foreign-language payload\,\tcode{T3.S1} or rewording\,\tcode{T7} (a dialect of one language is rewording) \\
\end{xwrows}\par
\begin{xwrows}
slang & \cb{3}\,slang or jargon\,\tcode{T7.S6} \\
context overload & \cb{3}\,overload\,\tcode{T5} or \cb{2}\,many-shot demonstrations\,\tcode{T4.S2} or \cb{1}\,worked examples\,\tcode{T5.S4} \\
n-shot & \cb{2}\,many-shot demonstrations\,\tcode{T4.S2} or \cb{1}\,worked examples\,\tcode{T5.S4} (labeled or unlabeled compliance examples) \\
repeated tokens, irrelevant distraction & \cb{3}\,overload\,\tcode{T5} or non-semantic sequence\,\tcode{T6} or \cb{1}\,repetition\,\tcode{T4.S3} (filler (S1) or distractors (S5)) \\
social hacking & \cb{4}\,social engineering or \cb{5}\,scenario framing \\
embedded conversation & \cb{2}\,fabricated dialogue\,\tcode{T4} (fictitious turns in which the model complies) \\
\end{xwrows}\par
\begin{xwrows}
historical context & \cb{5}\,era setting\,\tcode{T3.S5} \\
role-playing scenario & \cb{5}\,persona\,\tcode{T1} \\
leading response & \cb{6}\,required opening\,\tcode{T1.S1} \\
virtualization & \cb{5}\,persona\,\tcode{T1} or simulation\,\tcode{T2} or fiction and hypotheticals\,\tcode{T3} \\
hypothetical & \cb{5}\,hypothetical\,\tcode{T3.S2} \\
ignore-instructions & \cb{1}\,instruction nullification\,\tcode{T2.S1} \\
\end{xwrows}\par
\begin{xwrows}
detailed instruction & \cb{1}\,plain command\,\tcode{T1} (closest match, a detailed instruction set) \\
persuasion (40 sub-categories) & \cb{4}\,social engineering (the sub-categories are Zeng et al.'s) \\
specialized-tokens & \cb{3}\,non-semantic sequence\,\tcode{T6} (classified by surface form, not optimizer) \\
mixed technique & \cb{1}\,instruction manipulation or \cb{2}\,structural spoofing or \cb{3}\,obfuscation or \cb{4}\,social engineering or \cb{5}\,scenario framing or \cb{6}\,output coercion (two or more of the above combined) \\
\end{xwrows}\par
\end{xwcard}

\needspace{7\baselineskip}
\begin{xwcard}{NIST AI 100-2 2025~\cite{nist2025aml}}{organizes by objective \texttimes{} capability}
\xwnote{NIST's generative-AI taxonomy places attacks by attacker objective and required capability. Construction appears in its indirect-injection techniques and in a list of manual jailbreak approaches.}
\begin{xwrows}
availability breakdown, integrity violation, privacy compromise, misuse enablement & \outside{goal} \\
query access, resource control, direct prompting attacks, indirect prompt injection & \outside{delivery}, user queries versus manipulated resources \\
training data control, data poisoning, backdoor poisoning, targeted poisoning, model poisoning, supply chain attacks & \outside{lifecycle} \\
model control, fine-tuning circumvention, and white-box, gray-box and black-box attacks & \outside{model access} \\
training-time attacks, inference-time attacks & \outside{lifecycle} \\
jailbreaking, prompt extraction, training data extraction, model extraction, membership inference, misaligned outputs, increased computation & \outside{goal} \\
\end{xwrows}\par
\begin{xwrows}
optimization-based attacks, manual methods, automated model-based red teaming & \outside{search method}, optimized strings are placed by surface form \\
competing objectives, mismatched generalization & \outside{failure mode} \\
time-consuming background tasks, compromising connected resources, leaking information from user interactions & \outside{goal} \\
inhibiting capabilities & \cb{6}\,reasoning ban\,\tcode{T2.S6} or \cb{4}\,official notice\,\tcode{T1.S3} (the model is told an API is forbidden), \xwalso{goal} \\
disruptive output formatting & \cb{6}\,altered answer\,\tcode{T3.S3} or required opening\,\tcode{T1.S1} (homoglyphs, or an end-of-text token first), \xwalso{goal} \\
execution triggers & \cb{3}\,non-semantic sequence\,\tcode{T6} (trigger strings optimized by Neural Exec), \xwalso{search method} \\
\end{xwrows}\par
\begin{xwrows}
knowledge base poisoning & \outside{delivery channel, search method} \\
injection hiding & \cb{3}\,standard encoding\,\tcode{T2.S1} or \cb{1}\,fetch and follow\,\tcode{T6.S4} (hidden text alone is delivery), \xwalso{delivery rendering, lifecycle} \\
self-propagating injections & \outside{goal, lifecycle} \\
prefix injection & \cb{6}\,required opening\,\tcode{T1.S1} \\
refusal suppression & \cb{6}\,no-refusal order\,\tcode{T2.S4} \\
style injection & \cb{6}\,style mandate\,\tcode{T3.S5} \\
\end{xwrows}\par
\begin{xwrows}
role-play & \cb{5}\,persona\,\tcode{T1} \\
special encoding & \cb{3}\,standard encoding\,\tcode{T2.S1} \\
character transformation & \cb{3}\,classical cipher\,\tcode{T2.S2} or leetspeak\,\tcode{T1.S1} (ROT13 or Morse code, or leetspeak) \\
word transformation & \cb{3}\,classical cipher\,\tcode{T2.S2} or synonym swap\,\tcode{T7.S1} or fragmentation\,\tcode{T4} (Pig Latin, synonym swap, or payload splitting) \\
prompt-level transformation & \cb{3}\,foreign-language payload\,\tcode{T3.S1} \\
\xwex{the Crescendo attack} & \cb{4}\,commitment\,\tcode{T5} (in each turn), \xwalso{delivery turn structure, search method} \\
\end{xwrows}\par
\end{xwcard}

\needspace{7\baselineskip}
\begin{xwcard}{Gandalf 2025~\cite{pfister2025gandalf}}{organizes by trigger technique for one fixed goal}
\xwnote{Gandalf sorts player prompts by the trigger they use to reveal a fixed password, and all nine of its categories are constructions.}
\begin{xwrows}
direct & \cb{1}\,plain command\,\tcode{T1} \\
indirect & \cb{3}\,synonym swap\,\tcode{T7.S1} or euphemism\,\tcode{T7.S2} or allusion\,\tcode{T7.S3} or \cb{1}\,proxy request\,\tcode{T7} (another term, or a reworded question) \\
non-English input & \cb{3}\,language switching\,\tcode{T3} \\
context override & \cb{1}\,override\,\tcode{T2} \\
partial & \cb{1}\,proxy request\,\tcode{T7} or \cb{6}\,partial answer\,\tcode{T3.S4} (a hint, rhyme or definition, or one letter) \\
persuasion & \cb{4}\,authority claim\,\tcode{T1} or emotional appeal\,\tcode{T2} or urgency and threat\,\tcode{T3} or argument and obligation\,\tcode{T7} \\
\end{xwrows}\par
\begin{xwrows}
input obfuscation & \cb{3}\,encoding\,\tcode{T2} or character perturbation\,\tcode{T1} \\
output obfuscation & \cb{6}\,altered answer\,\tcode{T3.S3} \\
output fiction & \cb{5}\,story or script\,\tcode{T3.S1} \\
\xwex{trigger examples in its threat model (``ignore all instructions'', leetspeak, ciphers, role-playing, adversarial tokens)} & \cb{1}\,override\,\tcode{T2} or \cb{3}\,leetspeak\,\tcode{T1.S1} or classical cipher\,\tcode{T2.S2} or \cb{5}\,persona\,\tcode{T1} or \cb{3}\,non-semantic sequence\,\tcode{T6} \\
\end{xwrows}\par
\end{xwcard}

\needspace{7\baselineskip}
\begin{xwcard}{LLMail-Inject 2025~\cite{abdelnabi2025llmail}}{organizes by three annotator labels and team strategies}
\xwnote{LLMail-Inject labels submissions as direct instructions, obfuscation, or social engineering, and its team reports describe the construction strategies used against a simulated email assistant.}
\begin{xwrows}
direct instructions & \cb{1}\,plain command\,\tcode{T1} or direct aside\,\tcode{T5.S2} or (plain command\,\tcode{T1} \,+\, \cb{4}\,requester impersonation\,\tcode{T1.S4}) (commands addressed to the assistant) \\
obfuscation & \cb{3}\,foreign-language payload\,\tcode{T3.S1} or character perturbation\,\tcode{T1} or standard encoding\,\tcode{T2.S1} (other languages, inserted tokens, Base64) \\
social engineering & \cb{1}\,address to the reader\,\tcode{T5.S3} \,+\, \cb{4}\,social engineering (legitimate-looking mail to a human reader) \\
clear direct instructions, unclear, no injection & \outside{goal}, whether the email pursues the tool call \\
\xwex{special tokens that simulate turns, imitated end- and start-of-tool-output tokens, JSON escaping, user and summary tags} & (\cb{2}\,chat-template token\,\tcode{T2.S1} \,+\, user turn\,\tcode{T4.S4}) or boundary token\,\tcode{T2.S2} or section separator\,\tcode{T1.S2} or (role label\,\tcode{T1.S1} \,+\, user turn\,\tcode{T4.S4}) or markup element\,\tcode{T3.S2} \\
\xwex{fake-email blocks, synthetic conversations, misleading summaries, copied judge output format} & \cb{2}\,section separator\,\tcode{T1.S2} or (section separator\,\tcode{T1.S2} \,+\, \cb{3}\,filler padding\,\tcode{T5.S1}) or \cb{2}\,assistant turn\,\tcode{T4.S1} or (assistant turn\,\tcode{T4.S1} \,+\, \cb{1}\,task extension\,\tcode{T3.S1}) or \cb{2}\,status or error message\,\tcode{T5.S2} \\
\end{xwrows}\par
\begin{xwrows}
\xwex{character obfuscation, misspellings, multilingual and space-less languages, YAML and whitespace formatting} & \cb{3}\,character perturbation\,\tcode{T1} or language switching\,\tcode{T3} or surrogate syntax\,\tcode{T8} \\
\xwex{'after summarizing' prefaces and cascading instructions, declarative sentences, later text marked as dummy content} & \cb{1}\,task extension\,\tcode{T3.S1} or declarative statement\,\tcode{T5.S1} or instruction nullification\,\tcode{T2.S1} \\
\xwex{the example email in Figure 1} & \cb{1}\,address to the reader\,\tcode{T5.S3} \,+\, \cb{4}\,deadline\,\tcode{T3.S1} \,+\, sympathy\,\tcode{T2.S1} \,+\, incentive offer\,\tcode{T7.S7} \\
\xwex{subject-line and position choices} & \outside{delivery position} \\
\xwex{LLM-generated variants of a working template} & \outside{search method} \\
\xwex{exploits tailored to each defense} & \outside{target} \\
\end{xwrows}\par
\begin{xwrows}
\xwex{a reported attack that asks for the reply in bullet points with emojis} & \cb{6}\,forced format\,\tcode{T3} \\
\end{xwrows}\par
\end{xwcard}

\needspace{7\baselineskip}
\begin{xwcard}{Hong et al. 2026~\cite{hong2025sokeval}}{organizes by attacker access, then technique, and vulnerabilities}
\xwnote{Hong et al. classify jailbreaks by attacker access and then technique, and list nine model vulnerabilities. Construction appears in prompt modification and most vulnerability classes.}
\begin{xwrows}
obfuscation and encoding & \cb{3}\,character perturbation\,\tcode{T1} or encoding\,\tcode{T2} or language switching\,\tcode{T3} or surrogate syntax\,\tcode{T8} (code and markup reframing is F3.T8) \\
substitution and synonyms & \cb{3}\,synonym swap\,\tcode{T7.S1} or euphemism\,\tcode{T7.S2} or allusion\,\tcode{T7.S3} (synonyms, euphemisms, or code words) \\
decomposition & \cb{3}\,fragmentation\,\tcode{T4} \\
intent concealment & \cb{4}\,research purpose\,\tcode{T4.S1} or \cb{5}\,story or script\,\tcode{T3.S1} or hypothetical\,\tcode{T3.S2} or era setting\,\tcode{T3.S5} or \cb{3}\,distractors\,\tcode{T5.S5} (academic, narrative, distractor, or hypothetical framing) \\
heuristic algorithms, gradient estimation, reinforcement learning, LLM-generated attacks, multi-agent collaboration, surrogate models & \outside{search method} \\
multi-turn contextual attacks & \cb{3}\,fragmentation\,\tcode{T4} or \cb{4}\,commitment\,\tcode{T5} or \cb{2}\,assistant turn\,\tcode{T4.S1} (content of each turn), \xwalso{delivery turn structure} \\
\end{xwrows}\par
\begin{xwrows}
gradual escalation & \cb{4}\,commitment\,\tcode{T5} (in each turn), \xwalso{delivery turn structure} \\
behavior vulnerabilities, learning mechanism exploits, contextual ambiguity exploitation & \outside{failure mode} \\
overreliance on user instructions & \outside{failure mode, trust} \\
defensive mechanism bypass & \cb{3}\,foreign-language payload\,\tcode{T3.S1} (code applies to low-resource languages only), \xwalso{target, delivery channel, delivery modality} \\
decoding and sampling exploits, malicious fine-tuning, parameter modification, disabling safety mechanisms, backdoor insertion & \outside{model access} \\
prefix/suffix manipulation & \cb{3}\,non-semantic sequence\,\tcode{T6}, \xwalso{search method, delivery position} \\
\end{xwrows}\par
\begin{xwrows}
gradient-based optimization & \outside{search method, model access} \\
response format manipulation & \cb{6}\,forced format\,\tcode{T3} (structured output templates as covert channels) \\
summarization and translation exploitation & \cb{1}\,transformation object\,\tcode{T3.S3} or (transformation object\,\tcode{T3.S3} \,+\, \cb{3}\,foreign-language payload\,\tcode{T3.S1}) \\
psychological manipulation & \cb{4}\,emotional appeal\,\tcode{T2} or endorsement\,\tcode{T1.S7} or \cb{5}\,scenario framing (emotional appeals, authority endorsement, scenario framing) \\
hypothetical and scenario-based exploitation & \cb{5}\,fiction and hypotheticals\,\tcode{T3} \\
few-shot learning exploitation & \cb{2}\,many-shot demonstrations\,\tcode{T4.S2} or \cb{1}\,worked examples\,\tcode{T5.S4} (role-labeled turns or unlabeled examples) \\
\end{xwrows}\par
\begin{xwrows}
chain-of-thought exploitation & \cb{3}\,distractors\,\tcode{T5.S5} or \cb{2}\,assistant turn\,\tcode{T4.S1} or \cb{6}\,thought seeding\,\tcode{T1.S4} or \cb{3}\,fragmentation\,\tcode{T4} (distractors, pre-emptive answers, or decomposition) \\
conditional compliance exploitation & \cb{5}\,story or script\,\tcode{T3.S1} \,+\, \cb{4}\,authority claim\,\tcode{T1} (codes apply to nested authority scenarios only), \xwalso{model access} \\
personification and role-play exploitation & \cb{5}\,persona\,\tcode{T1} or (persona\,\tcode{T1} \,+\, story or script\,\tcode{T3.S1}) \\
function call exploitation & \outside{delivery channel} \\
system prompt leakage & \outside{goal} \\
context length exploitation & \cb{3}\,filler padding\,\tcode{T5.S1} or \cb{2}\,many-shot demonstrations\,\tcode{T4.S2} \\
\end{xwrows}\par
\begin{xwrows}
\xwex{CodeAttack and CodeChameleon, which Table 15 files under obfuscation and encoding} & (\cb{3}\,pseudocode\,\tcode{T8.S2} \,+\, positional extraction\,\tcode{T4.S2}) or (pseudocode\,\tcode{T8.S2} \,+\, encoding\,\tcode{T2}) \\
\end{xwrows}\par
\end{xwcard}

\needspace{7\baselineskip}
\begin{xwcard}{HiddenLayer APE 2026~\cite{hiddenlayer2025ape}}{organizes by tactic (the mechanism exploited), impact, objective}
\xwnote{APE groups prompt techniques into eight tactics by the mechanism they exploit and keeps impacts and objectives separate, so construction appears throughout its tactics.}
\begin{xwrows}
obfuscation, encrypted / encoded input, random character insertion / replacement, payload splitting, translated language & \cb{3}\,encoding\,\tcode{T2} or character perturbation\,\tcode{T1} or fragmentation\,\tcode{T4} or language switching\,\tcode{T3} (by member technique) \\
encrypted / encoded output & \cb{6}\,altered answer\,\tcode{T3.S3} (the answer, not the input, is encoded) \\
lore abuse, KROP, synonym attack & \cb{3}\,allusion\,\tcode{T7.S3} or synonym swap\,\tcode{T7.S1} (by member technique) \\
language completion games & \cb{5}\,game objective\,\tcode{T4.S1} or \cb{1}\,proxy request\,\tcode{T7} or \cb{6}\,pattern continuation\,\tcode{T1.S3} (by game type) \\
cognitive manipulation & \cb{1}\,instruction manipulation or \cb{3}\,obfuscation or \cb{4}\,social engineering or \cb{5}\,scenario framing or \cb{6}\,output coercion (by member technique) \\
imperative emphasis, instruction override, meta prompting & \cb{1}\,emphasis\,\tcode{T4} or override\,\tcode{T2} or execution rules\,\tcode{T6} (by member technique) \\
\end{xwrows}\par
\begin{xwrows}
response priming, refusal suppression & \cb{6}\,response priming\,\tcode{T1} or refusal suppression\,\tcode{T2} \\
emotional appeal & \cb{4}\,emotional appeal\,\tcode{T2} or urgency and threat\,\tcode{T3} (distress, flattery, urgency or threats) \\
role playing & \cb{5}\,persona\,\tcode{T1} \\
narrative smuggling, cognitive overload & \cb{3}\,positional extraction\,\tcode{T4.S2} or filler padding\,\tcode{T5.S1} or cognitive overload\,\tcode{T5.S2} (by member technique) \\
syntax-based input & \cb{3}\,surrogate syntax\,\tcode{T8} or \cb{2}\,structured-data disguise\,\tcode{T3} (F2 when posing as configuration) \\
semantic confusion & \cb{1}\,contradiction\,\tcode{T2.S6} or \cb{3}\,paradox\,\tcode{T5.S4} (conflicting directives or paradox) \\
\end{xwrows}\par
\begin{xwrows}
crescendo & \cb{4}\,commitment\,\tcode{T5} (in each turn), \xwalso{delivery turn structure} \\
embedding / token manipulation, token-aware prompting, glitch tokens, high-perplexity padding & \cb{3}\,spacing changes\,\tcode{T1.S2} or non-semantic sequence\,\tcode{T6} or filler padding\,\tcode{T5.S1} (by member technique) \\
algorithmic attacks & \cb{3}\,non-semantic sequence\,\tcode{T6} (the optimized string itself), \xwalso{search method} \\
context manipulation, control token injection / spoofing, chain-of-thought spoofing, tool spoofing, conversation spoofing, system prompt spoofing, policy puppetry & \cb{2}\,control tokens\,\tcode{T2} or delimiters\,\tcode{T1} or assistant turn\,\tcode{T4.S1} or counterfeit tool results\,\tcode{T5} or fabricated dialogue\,\tcode{T4} or role label\,\tcode{T1.S1} or runtime config\,\tcode{T3.S1} (spoofed context, other members below) \\
few-shot prompting, context window separation & \cb{1}\,worked examples\,\tcode{T5.S4} or \cb{2}\,many-shot demonstrations\,\tcode{T4.S2} or section separator\,\tcode{T1.S2} or \cb{1}\,reset\,\tcode{T2.S2} (by member and role labels) \\
defined dictionary & \cb{1}\,meaning remap\,\tcode{T6.S5} or \cb{3}\,variable concatenation\,\tcode{T4.S1} \\
\end{xwrows}\par
\begin{xwrows}
refusal hijacking & \cb{1}\,conditional trigger\,\tcode{T6.S1} or \cb{6}\,staged denial\,\tcode{T1.S5} (payload fires when the model refuses) \\
pretexting & \cb{4}\,pretext\,\tcode{T4} \\
output structuring, summarization, output natural language, syntax-based output, output termination, output pruning & \cb{6}\,altered answer\,\tcode{T3.S3} or partial answer\,\tcode{T3.S4} or schema only\,\tcode{T3.S1} or refusal suppression\,\tcode{T2} or forced format\,\tcode{T3} (by member technique) \\
output creative & \cb{6}\,style mandate\,\tcode{T3.S5} or \cb{5}\,story or script\,\tcode{T3.S1} \\
stop-token prevention & \cb{6}\,length mandate\,\tcode{T3.S6} (unbounded output), \xwalso{goal} \\
attack augmentation & \outside{delivery turn structure, delivery channel}, amplifies other techniques \\
\end{xwrows}\par
\begin{xwrows}
attack stacking & \cb{1}\,repetition\,\tcode{T4.S3} \\
multi-turn & \cb{4}\,commitment\,\tcode{T5} or \cb{3}\,fragmentation\,\tcode{T4} (per turn, commitment or fragments), \xwalso{delivery turn structure} \\
indirect visibility & \outside{delivery channel}, retrieval ranking \\
multi-LLM attacks, recursive injection & \outside{delivery channel, lifecycle}, propagates between chained models \\
safety / judge model manipulation & \outside{target} \\
privacy compromise / data exposure, integrity violation / behavior subversion, availability breakdown / operational disruption (with their 22 objectives) & \outside{goal} \\
\end{xwrows}\par
\end{xwcard}

\needspace{7\baselineskip}
\begin{xwcard}{Xu et al. 2026~\cite{xu2026sokjailbreak}}{organizes by dominant attack mechanism (seven types)}
\xwnote{Xu et al. file each jailbreak under one of seven mechanism types. Construction defines the shuffle-based type and appears in members of the strategy- and flaw-based types.}
\begin{xwrows}
logprob-based attacks & \outside{search method, model access} \\
LLM-based attacks, template-based attacks & \outside{search method} \\
shuffle-based attacks & \cb{3}\,transposition\,\tcode{T2.S5} or spacing changes\,\tcode{T1.S2} or character perturbation\,\tcode{T1} (reordered tokens, separators, whitespace, surface noise) \\
multi-round-based attacks & \cb{4}\,commitment\,\tcode{T5} or \cb{3}\,fragmentation\,\tcode{T4} (content of each turn), \xwalso{delivery turn structure} \\
flaw-based attacks & \outside{failure mode}, defined by model-dependent weaknesses \\
strategy-based attacks & \cb{4}\,social engineering or \cb{5}\,persona\,\tcode{T1} (persuasive framing or persona appeals) \\
\end{xwrows}\par
\begin{xwrows}
\xwex{PAP, DAN (strategy-based members)} & \cb{4}\,social engineering or \cb{5}\,unrestricted AI\,\tcode{T1.S1} \\
\xwex{CodeAttacker, QueryAttack (strategy-based members)} & (\cb{3}\,pseudocode\,\tcode{T8.S2} \,+\, positional extraction\,\tcode{T4.S2}) or surrogate syntax\,\tcode{T8} \\
\xwex{DrAttack (strategy-based member)} & \cb{3}\,fragmentation\,\tcode{T4} \\
\xwex{ReNeLLM (strategy-based member)} & (\cb{3}\,rewording\,\tcode{T7} \,+\, pseudocode\,\tcode{T8.S2}) or (rewording\,\tcode{T7} \,+\, \cb{5}\,story or script\,\tcode{T3.S1}) or (\cb{3}\,rewording\,\tcode{T7} \,+\, \cb{6}\,schema only\,\tcode{T3.S1}) (rewriting plus code, story, or table nesting) \\
\xwex{CipherChat, Multijail, ArtPrompt (flaw-based members)} & \cb{3}\,classical cipher\,\tcode{T2.S2} or foreign-language payload\,\tcode{T3.S1} or ASCII art\,\tcode{T1.S5} \\
\xwex{GCG (logprob-based member)} & \cb{3}\,non-semantic sequence\,\tcode{T6}, \xwalso{search method} \\
\end{xwrows}\par
\begin{xwrows}
\xwex{DeepInception (multi-round member)} & \cb{5}\,story or script\,\tcode{T3.S1}, \xwalso{delivery turn structure} \\
\xwex{task-diverting mechanisms described for PAP and ReNeLLM} & \cb{1}\,proxy request\,\tcode{T7} \\
\end{xwrows}\par
\end{xwcard}

\needspace{7\baselineskip}
\begin{xwcard}{Wang et al. 2026~\cite{wang2026landscape}}{organizes by payload generation, facets, and AgentPI attacks}
\xwnote{Wang et al. classify injection papers by payload generation with five facets. Construction appears in structural encoding, three visibility values, and the five AgentPI attacks.}
\begin{xwrows}
manual template & \outside{search method} \\
LLM generation, gradient, genetic, sampling & \outside{search method, model access} \\
structural encoding & \cb{3}\,standard encoding\,\tcode{T2.S1} or ASCII art\,\tcode{T1.S5} or surrogate syntax\,\tcode{T8} (file layouts may be rendering tricks) \\
direct, indirect, and supply-chain prompt injection & \outside{delivery channel} \\
LLM, LLM-integrated application, LLM agent & \outside{target} \\
integrity, confidentiality, availability & \outside{goal} \\
\end{xwrows}\par
\begin{xwrows}
white-box and black-box capabilities & \outside{model access}, access to gradients, logits, or outputs \\
visible payloads & \cb{1}\,override\,\tcode{T2} or importance flag\,\tcode{T4.S2} (explicit triggers such as ignore previous or important) \\
invisible via encoding & \cb{3}\,encoding\,\tcode{T2} \\
invisible via semantics & \cb{3}\,non-semantic sequence\,\tcode{T6} (optimized non-meaningful suffixes), \xwalso{search method} \\
invisible via merging in context & \outside{delivery channel}, hidden by where it sits (poisoned sources, logs) \\
invisible in vision & \outside{delivery modality} \\
\end{xwrows}\par
\begin{xwrows}
action switching & \cb{1}\,override\,\tcode{T2} or replacement command\,\tcode{T1.S2} (payload overrides the user's tool choice), \xwalso{goal} \\
parameter manipulation & \cb{4}\,false premise\,\tcode{T7.S8} or (false premise\,\tcode{T7.S8} \,+\, \cb{1}\,plain command\,\tcode{T1}) or replacement command\,\tcode{T1.S2} (false data alters a tool parameter), \xwalso{goal} \\
branch divergence & \cb{4}\,false premise\,\tcode{T7.S8} or (false premise\,\tcode{T7.S8} \,+\, \cb{2}\,status or error message\,\tcode{T5.S2}) (fabricated facts flip a conditional branch), \xwalso{goal} \\
reasoning corruption & \cb{4}\,false premise\,\tcode{T7.S8} or \cb{1}\,meaning remap\,\tcode{T6.S5} or (override\,\tcode{T2} \,+\, \cb{4}\,false premise\,\tcode{T7.S8}) (false conclusions or inverted min/max logic), \xwalso{goal} \\
delegation exploitation & \cb{1}\,plain command\,\tcode{T1} or (\cb{2}\,status or error message\,\tcode{T5.S2} \,+\, \cb{1}\,override\,\tcode{T2} \,+\, \cb{4}\,harm-prevention purpose\,\tcode{T4.S3}) (user delegated authority to the content), \xwalso{trust} \\
\xwex{the manual-template description (direct overrides, authoritative-sounding instructions)} & \cb{1}\,override\,\tcode{T2} or \cb{4}\,authority claim\,\tcode{T1} \\
\end{xwrows}\par
\end{xwcard}

\needspace{7\baselineskip}
\begin{xwcard}{Maloyan \& Namiot 2026~\cite{maloyan2026coding}}{organizes by delivery vector, modality, and propagation}
\xwnote{Maloyan and Namiot classify coding-assistant attacks by delivery vector, modality, and propagation. Construction fills the direct-injection sub-classes and the text and semantic modalities.}
\begin{xwrows}
D1.1 role hijacking & \cb{4}\,authority claim\,\tcode{T1} or \cb{2}\,role label\,\tcode{T1.S1} (privilege claim in prose or as label), \xwalso{delivery channel} \\
D1.2 context override & \cb{1}\,override\,\tcode{T2} or replacement command\,\tcode{T1.S2} (redefines the agent's purpose), \xwalso{delivery channel} \\
D1.3 instruction negation & \cb{1}\,instruction nullification\,\tcode{T2.S1}, \xwalso{delivery channel} \\
D2.1 repository-based (rules file backdoor, code comments, issue/PR poisoning), D2.2 documentation-based (README exploitation, API doc poisoning, manifest injection), D2.3 web content (search poisoning, documentation compromise) & \outside{delivery channel} \\
D3.1 MCP attacks (tool poisoning, rug pull, shadowing, tool squatting) & \outside{delivery channel, lifecycle}, rug pull changes tools after approval \\
D3.2 transport attacks (MITM, DNS rebinding, SSE injection) & \outside{delivery channel, target} \\
\end{xwrows}\par
\begin{xwrows}
M1.1 hierarchy exploitation & \cb{4}\,authority claim\,\tcode{T1} or \cb{2}\,structural spoofing (privilege level claims) \\
M1.2 completion attacks & \cb{2}\,assistant turn\,\tcode{T4.S1} or \cb{6}\,response priming\,\tcode{T1} \\
M1.3 encoding obfuscation & \cb{3}\,standard encoding\,\tcode{T2.S1} or character perturbation\,\tcode{T1} or fragmentation\,\tcode{T4} (Base64, Unicode, word splitting) \\
M2.1 XOXO (cross-origin context poisoning) & \outside{delivery channel} \\
M2.2 implicit instructions & \cb{1}\,indirect phrasing\,\tcode{T5} \\
M2.3 logic bombs & \cb{1}\,conditional trigger\,\tcode{T6.S1} \,+\, \cb{3}\,pseudocode\,\tcode{T8.S2} (trigger hidden in benign-looking code) \\
\end{xwrows}\par
\begin{xwrows}
M3.1 image injection, M3.2 audio attacks, M3.3 video frames & \outside{delivery modality} \\
P1 single-shot, P2.1 config modification, P2.2 memory poisoning, P2.3 system backdoors, P3.1 repository worms, P3.2 dependency chain, P3.3 agent-to-agent & \outside{lifecycle} \\
data exfiltration, code injection, privilege escalation, denial of service, persistence & \outside{goal}, the five attack objectives \\
\xwex{user-agent, agent-tool, tool-tool, and session boundaries} & \outside{trust} \\
\xwex{content injector, tool publisher, network attacker} & \outside{delivery channel}, attacker capability levels \\
\xwex{the four components of Log-To-Leak (trigger, tool binding, justification, pressure)} & \cb{1}\,conditional trigger\,\tcode{T6.S1} \,+\, \cb{4}\,pretext\,\tcode{T4} \,+\, urgency and threat\,\tcode{T3} (tool binding is the goal), \xwalso{goal} \\
\end{xwrows}\par
\begin{xwrows}
\xwex{case-study payloads (hidden system comment, tool description with an importance flag, rules-file routine step)} & (\cb{2}\,role label\,\tcode{T1.S1} \,+\, \cb{1}\,instruction nullification\,\tcode{T2.S1}) or (importance flag\,\tcode{T4.S2} \,+\, routine step\,\tcode{T3.S2}) or routine step\,\tcode{T3.S2} \\
\end{xwrows}\par
\end{xwcard}

\needspace{7\baselineskip}
\begin{xwcard}{Promptware kill chain 2026~\cite{promptware2026killchain}}{organizes by kill-chain stage (seven stages)}
\xwnote{The promptware kill chain orders attacks into seven stages. Construction appears in its jailbreaking labels and in the techniques it describes for the first two stages.}
\begin{xwrows}
reconnaissance, command and control & \outside{lifecycle} \\
direct and indirect prompt injection (initial access) & \outside{delivery channel} \\
multimodal LLM injection (initial access) & \outside{delivery modality} \\
DAN, role-playing (privilege escalation) & \cb{5}\,unrestricted AI\,\tcode{T1.S1} or persona\,\tcode{T1} \\
instr. override & \cb{1}\,override\,\tcode{T2} (Table II privilege-escalation label) \\
role-play JB & \cb{5}\,persona\,\tcode{T1} (Table II privilege-escalation label) \\
\end{xwrows}\par
\begin{xwrows}
instr. obfusc. & \cb{3}\,obfuscation (Table II privilege-escalation label) \\
semantic bypass & \cb{3}\,ASCII art\,\tcode{T1.S5} (ArtPrompt's ASCII-art masking) \\
delayed tool inv & \cb{1}\,conditional trigger\,\tcode{T6.S1} (Table II privilege-escalation label) \\
tool confusion & \cb{1}\,meaning remap\,\tcode{T6.S5} or \cb{4}\,insider identity\,\tcode{T1.S1} (Freysa heist, a label not defined in the text) \\
control bypass, approval bypass, auto tool inv. & \outside{target} \\
auto RAG mix, zero-click RAG & \outside{delivery channel} \\
\end{xwrows}\par
\begin{xwrows}
retrieval-dependent and retrieval-independent persistence, git repo, session & \outside{lifecycle} \\
on-device and off-device lateral movement, permission-based, self-replication, pipeline, supply chain, git propagation & \outside{lifecycle} \\
RCE, data exfiltration (actions on objective) & \outside{goal} \\
browser/search, enterprise, coding assistant, AI agent, agentic/CUA, crypto/DeFi, AI worm, multimodal & \outside{target, lifecycle, delivery modality}, Table II incident categories \\
\xwex{evasion techniques described for initial access (ASCII smuggling, Base64, ASCII art, delayed tool invocation)} & \cb{3}\,invisible characters\,\tcode{T1.S4} or standard encoding\,\tcode{T2.S1} or ASCII art\,\tcode{T1.S5} or \cb{1}\,conditional trigger\,\tcode{T6.S1} \\
\xwex{text colored like the background (initial access)} & \outside{delivery rendering} \\
\end{xwrows}\par
\begin{xwrows}
\xwex{jailbreaking techniques described for privilege escalation (overrides, persona and role-play, adversarial suffixes, ASCII-art prompts)} & \cb{1}\,override\,\tcode{T2} or \cb{5}\,persona\,\tcode{T1} or fiction and hypotheticals\,\tcode{T3} or \cb{3}\,non-semantic sequence\,\tcode{T6} or ASCII art\,\tcode{T1.S5} \\
\xwex{multi-turn jailbreaking such as Crescendo} & \cb{4}\,commitment\,\tcode{T5} (in each turn), \xwalso{delivery turn structure} \\
\xwex{appendix payloads (calendar invitation, memory implant)} & (\cb{5}\,job persona\,\tcode{T1.S6} \,+\, \cb{1}\,conditional trigger\,\tcode{T6.S1}) or (conditional trigger\,\tcode{T6.S1} \,+\, fetch and follow\,\tcode{T6.S4}) or \cb{6}\,forced format\,\tcode{T3} \\
ASCII art encoding & \cb{3}\,ASCII art\,\tcode{T1.S5} \\
\end{xwrows}\par
\end{xwcard}

\needspace{7\baselineskip}
\begin{xwcard}{Giarrusso et al. 2026~\cite{giarrusso2026guardrails}}{organizes by mechanism by which safety is bypassed}
\xwnote{It groups jailbreak techniques into seven families by the mechanism that bypasses safety. Most of its 31 subcategories are construction, and a few are turn structure.}
\begin{xwrows}
impersonation attacks \& fictional scenarios & \cb{5}\,scenario framing or \cb{4}\,pretext\,\tcode{T4} \\
role play, defined personas, virtual AI, antagonistic entities split, jailbroken model simulation & \cb{5}\,persona\,\tcode{T1} (virtual AI S1/S2, split S4, jailbroken S1) \\
benign context framing, research \& testing, joking pretext, game & \cb{4}\,pretext\,\tcode{T4} or \cb{5}\,game objective\,\tcode{T4.S1} (the game variant is F5.T4.S1) \\
fictional framing & \cb{5}\,fiction and hypotheticals\,\tcode{T3} \\
privilege escalation & \cb{5}\,special mode\,\tcode{T1.S2} or unrestricted AI\,\tcode{T1.S1} or \cb{1}\,capitals and symbols\,\tcode{T4.S1} or \cb{4}\,insider identity\,\tcode{T1.S1} (admin claims, jailbroken mode, or capitals) \\
sudo/admin mode & \cb{5}\,special mode\,\tcode{T1.S2} or \cb{1}\,conditional trigger\,\tcode{T6.S1} (variant with a 'special' instruction) \\
\end{xwrows}\par
\begin{xwrows}
typographical authority simulation, instruction repetition & \cb{1}\,emphasis\,\tcode{T4} (capitals (S1) or repetition (S3)) \\
persuasion & \cb{4}\,social engineering or \cb{1}\,repetition\,\tcode{T4.S3} (instruction repetition is F1.T4.S3) \\
logical, evidential, and quantification-based persuasion & \cb{4}\,scripted deduction\,\tcode{T7.S9} or cited evidence\,\tcode{T7.S10} \\
authority and norm-based persuasion & \cb{4}\,endorsement\,\tcode{T1.S7} or normality assertion\,\tcode{T6.S2} \\
emotional, reciprocity-based, and commitment-based persuasion & \cb{4}\,emotional appeal\,\tcode{T2} or reciprocity\,\tcode{T7.S6} or commitment\,\tcode{T5} or hypocrisy charge\,\tcode{T7.S3} (repeated-request variant with a hypocrisy charge) \\
urgency and scarcity-based persuasion, manipulative and coercive persuasion & \cb{4}\,urgency and threat\,\tcode{T3} (coercion by consequences is S4) \\
\end{xwrows}\par
\begin{xwrows}
cognitive overload \& attention misalignment & \cb{3}\,overload\,\tcode{T5} or fragmentation\,\tcode{T4} or \cb{1}\,proxy request\,\tcode{T7} \\
distractor instructions, context saturation, mathematical \& decomposition attacks & \cb{3}\,overload\,\tcode{T5} or fragmentation\,\tcode{T4} (distractors S5, filler S1, puzzles S2) \\
indirect task deflection & \cb{1}\,composition request\,\tcode{T7.S2} \\
encoding \& obfuscation & \cb{3}\,obfuscation \\
surface obfuscation, token splitting, ASCII-art & \cb{3}\,character perturbation\,\tcode{T1} (spacing S2, ASCII art S5) \\
semantic rewriting, sentence-level, token-level & \cb{3}\,rewording\,\tcode{T7} \\
\end{xwrows}\par
\begin{xwrows}
linguistic encoding, low-resource languages, base64 & \cb{3}\,encoding\,\tcode{T2} or language switching\,\tcode{T3} or character perturbation\,\tcode{T1} (low-resource T3.S1, Base64 T2.S1) \\
lexical techniques & \cb{3}\,non-semantic sequence\,\tcode{T6} (closest match, as readable triggers vary) \\
embedded prompting & \cb{3}\,surrogate syntax\,\tcode{T8} or \cb{2}\,structured-data disguise\,\tcode{T3} or \cb{3}\,ASCII art\,\tcode{T1.S5} (comments and images are delivery), \xwalso{delivery position, delivery modality} \\
goal-conflicting attacks & \cb{1}\,instruction manipulation or \cb{6}\,output coercion \\
prefix injection & \cb{6}\,required opening\,\tcode{T1.S1} or \cb{2}\,fabricated dialogue\,\tcode{T4} (answer opening, or fake history) \\
instruction masking, text-completion-as-instruction & \cb{1}\,transformation object\,\tcode{T3.S3} or \cb{6}\,seeded continuation\,\tcode{T1.S2} (TCINS is the continuation case) \\
\end{xwrows}\par
\begin{xwrows}
refusal suppression & \cb{6}\,refusal suppression\,\tcode{T2} \\
context ignoring & \cb{1}\,instruction nullification\,\tcode{T2.S1} or policy revocation\,\tcode{T2.S4} \\
assumption of responsibility & \cb{4}\,argument and obligation\,\tcode{T7} or \cb{1}\,policy revocation\,\tcode{T2.S4} (closest are an autonomy argument or rule revocation) \\
objective juxtaposition & \cb{1}\,added command\,\tcode{T1.S1} or contradiction\,\tcode{T2.S6} or task extension\,\tcode{T3.S1} \\
data poisoning attacks & \cb{2}\,many-shot demonstrations\,\tcode{T4.S2} or \cb{4}\,false premise\,\tcode{T7.S8} or commitment\,\tcode{T5} (escalation across turns), \xwalso{delivery turn structure} \\
incremental poisoning & \cb{4}\,foot-in-the-door\,\tcode{T5.S1} or \cb{3}\,fragmentation\,\tcode{T4} (escalation or split across turns), \xwalso{delivery turn structure} \\
\end{xwrows}\par
\begin{xwrows}
many-shot jailbreaking & \cb{1}\,worked examples\,\tcode{T5.S4} or \cb{2}\,many-shot demonstrations\,\tcode{T4.S2} (its example is unlabeled Q/A pairs) \\
false fact/bias instillation & \cb{4}\,false premise\,\tcode{T7.S8} \\
automated attacks (feedback-driven optimization, genetic algorithms, LLM-based agentic refinement) & \outside{search method, model access} \\
automatic attacks & \outside{search method} \\
no technique & \cb{1}\,plain command\,\tcode{T1} \\
\end{xwrows}\par
\end{xwcard}

\needspace{7\baselineskip}
\begin{xwcard}{MLCommons 2026~\cite{mlcommons2026jailbreak}}{organizes by mechanism of a single-turn jailbreak}
\xwnote{MLCommons sorts single-turn jailbreak prompts by manipulation mechanism into four families, and construction appears in every family, mixed with automated search methods.}
\begin{xwrows}
perturbation & \cb{3}\,character perturbation\,\tcode{T1} or rewording\,\tcode{T7} or non-semantic sequence\,\tcode{T6} (surface form changes, meaning kept) \\
plain perturbations, character-level micro-edits, local paraphrase \& rewording & \cb{3}\,character perturbation\,\tcode{T1} or rewording\,\tcode{T7} (by leaf) \\
transfer text perturbations, adversarial triggers, paraphrase-based transfer & \cb{3}\,non-semantic sequence\,\tcode{T6} or rewording\,\tcode{T7} (split by production method), \xwalso{search method} \\
encoding abuse & \cb{3}\,encoding\,\tcode{T2} or character perturbation\,\tcode{T1} or \cb{2}\,structured-data disguise\,\tcode{T3} or \cb{3}\,surrogate syntax\,\tcode{T8} or filler padding\,\tcode{T5.S1} (by category) \\
encoding \& Unicode tricks, ASCII encoding, Base64 / URL / obfuscation, Unicode / Bidi / zero-width & \cb{3}\,ASCII art\,\tcode{T1.S5} or spacing changes\,\tcode{T1.S2} or encoding\,\tcode{T2} or invisible characters\,\tcode{T1.S4} (by leaf) \\
wrappers \& schemas & \cb{2}\,structured-data disguise\,\tcode{T3} or \cb{3}\,surrogate syntax\,\tcode{T8} or filler padding\,\tcode{T5.S1} (by leaf) \\
\end{xwrows}\par
\begin{xwrows}
JSON / Markdown / code-blocks & \cb{2}\,structured-data disguise\,\tcode{T3} or \cb{3}\,surrogate syntax\,\tcode{T8} (F2 when framed as configuration) \\
length / prefix-suffix wrappers & \cb{3}\,filler padding\,\tcode{T5.S1} (benign padding, payload placed last), \xwalso{delivery position} \\
overt carriers & \cb{1}\,override\,\tcode{T2} or \cb{5}\,persona\,\tcode{T1} or \cb{4}\,authority claim\,\tcode{T1} or pretext\,\tcode{T4} (by category) \\
direct override patterns & \cb{1}\,override\,\tcode{T2} \\
simple overrides & \cb{1}\,override\,\tcode{T2} or \cb{5}\,special mode\,\tcode{T1.S2} or \cb{4}\,authority claim\,\tcode{T1} (ignore-instructions, developer mode or claimed privilege) \\
DAN-style composites & \cb{5}\,persona\,\tcode{T1} \,+\, \cb{1}\,override\,\tcode{T2} (persona with its own rule system) \\
\end{xwrows}\par
\begin{xwrows}
role-play \& template jailbreaks & \cb{4}\,pretext\,\tcode{T4} or \cb{5}\,persona\,\tcode{T1} (by leaf) \\
benign pretext / safety reframing & \cb{4}\,pretext\,\tcode{T4} \\
persona role-play & \cb{5}\,persona\,\tcode{T1} \\
composition \& ordering & \cb{5}\,fiction and hypotheticals\,\tcode{T3} or \cb{1}\,task embedding\,\tcode{T3} or \cb{4}\,commitment\,\tcode{T5} or \cb{3}\,fragmentation\,\tcode{T4} (by category) \\
context framing \& deception & \cb{5}\,fiction and hypotheticals\,\tcode{T3} or \cb{1}\,task embedding\,\tcode{T3} or \cb{4}\,document guise\,\tcode{T1.S5} or commitment\,\tcode{T5} (by leaf) \\
benign wrapper + harmful core & \cb{5}\,story or script\,\tcode{T3.S1} or \cb{4}\,document guise\,\tcode{T1.S5} or \cb{1}\,task embedding\,\tcode{T3} (story, document guise or nested task) \\
\end{xwrows}\par
\begin{xwrows}
scenario-based assembly & \cb{4}\,commitment\,\tcode{T5} or \cb{5}\,fiction and hypotheticals\,\tcode{T3} (stepwise escalation within one prompt) \\
fragment assembly, token/phrase shuffling, interleaving benign / harmful sections & \cb{3}\,fragmentation\,\tcode{T4} or transposition\,\tcode{T2.S5} or filler padding\,\tcode{T5.S1} (by leaf) \\
jailbreak prompt optimization \& search & \outside{search method} \\
\xwex{AdaPPA, a prefill attack listed under length / prefix-suffix wrappers} & \cb{6}\,seeded continuation\,\tcode{T1.S2} or \cb{2}\,assistant turn\,\tcode{T4.S1} (response prefill, the closest matches) \\
\end{xwrows}\par
\end{xwcard}

\needspace{7\baselineskip}
\begin{xwcard}{Iyer et al. 2026~\cite{iyer2026talk}}{organizes by target crossed with the defense layer bypassed}
\xwnote{Iyer et al. cross four targets with six techniques defined by the defense layer an attack bypasses. Three technique columns correspond to construction classes.}
\begin{xwrows}
safety \& alignment bypass, system \& tool hijacking, information exfiltration, service disruption & \outside{goal} \\
instructional & \cb{1}\,instruction manipulation (bypasses instruction priority) \\
persuasion \& deception & \cb{4}\,social engineering or \cb{5}\,scenario framing (persuasion or persona framing) \\
obfuscation & \cb{3}\,obfuscation (includes optimized suffixes) \\
cross-modal & \outside{delivery modality} \\
indirect injection & \outside{delivery channel} \\
\end{xwrows}\par
\begin{xwrows}
model internals & \outside{model access} \\
jailbreaking, information disclosure, system prompt extraction, sensitive data exfiltration, training data recall, probing \& reconnaissance, model/architecture extraction, service disruption, denial of service, functional degradation, watermarking attacks & \outside{goal} \\
instruction hijacking & \cb{1}\,instruction manipulation (by name, as its leaves are not listed) \\
deceptive/psychological manipulation & \cb{4}\,social engineering or \cb{5}\,scenario framing (includes multi-turn leaves), \xwalso{delivery turn structure} \\
obfuscation \& evasion & \cb{3}\,obfuscation \\
multi-modal jailbreaking, multi-modal indirect injection & \outside{delivery modality} \\
\end{xwrows}\par
\begin{xwrows}
in-context learning exploitation & \cb{1}\,worked examples\,\tcode{T5.S4} or \cb{2}\,many-shot demonstrations\,\tcode{T4.S2} \\
hybrid \& composite attacks, composite payload injection & \cb{1}\,instruction manipulation or \cb{3}\,obfuscation or \cb{4}\,social engineering or \cb{5}\,scenario framing (combinations of other subcategories) \\
system/infrastructure exploitation & \outside{target}, reasoning, memory, context, mode parameters \\
unauthorized execution \& agency abuse & \outside{goal, target} \\
prompt injection, data source injection vectors, agentic \& tool-based injection & \outside{delivery channel} \\
time-based \& stateful injection & \cb{1}\,conditional trigger\,\tcode{T6.S1}, \xwalso{lifecycle} \\
\end{xwrows}\par
\begin{xwrows}
decoding \& generation manipulation, controllable decoding, enforced decoding, refusal subspace ablation, MoE routing manipulation, LLM router manipulation, activation steering & \outside{model access} \\
attack generation \& optimization methods, gradient-based, automated generative / LLM-as-attacker, sampling \& fuzzing, black-box optimization, text-to-image and retrieval-based methods, manual/human-in-the-loop methods & \outside{search method} \\
\xwex{leaves named in its Service Disruption x Instructional walkthrough, e.g. excessive-token-generation-requests} & \cb{6}\,length mandate\,\tcode{T3.S6} or \cb{1}\,contradiction\,\tcode{T2.S6} or \cb{3}\,overload\,\tcode{T5}, \xwalso{goal} \\
\end{xwrows}\par
\end{xwcard}

\needspace{7\baselineskip}
\begin{xwcard}{OWASP LLM01:2026~\cite{owasp2026llm01}}{organizes by delivery surface, propagation and encoding}
\xwnote{OWASP LLM01:2026 describes each injection by delivery surface, propagation and encoding, so construction appears only through its encoding axis and a few examples.}
\begin{xwrows}
direct prompt injection, indirect prompt injection & \outside{delivery channel} \\
jailbreaking & \outside{goal}, injection aimed at safety violations \\
untrusted surfaces, semi-trusted surfaces, trusted surfaces & \outside{trust, delivery channel} \\
delivery surface, direct input, retrieved content, tool output, tool connection channel, persistent memory & \outside{delivery channel} \\
propagation behavior, single-shot, multi-step kill-chain, cross-session through memory or RAG, self-replicating across agents & \outside{lifecycle} \\
encoding, base64 or other obfuscation, invisible Unicode, low-resource language & \cb{3}\,encoding\,\tcode{T2} or \cb{3}\,obfuscation or invisible characters\,\tcode{T1.S4} or foreign-language payload\,\tcode{T3.S1} (text-level values of the axis) \\
\end{xwrows}\par
\begin{xwrows}
plain text, multimodal or steganographic & \outside{delivery modality}, text versus pixels or audio \\
\xwex{listed outcomes (disclosure, output manipulation, tool invocation, exfiltration, persistent compromise)} & \outside{goal} \\
\xwex{common example 1 (direct prompt-input override), scenario 1 (direct injection)} & \cb{1}\,override\,\tcode{T2} \\
\xwex{common examples 5 and 8 (invisible characters and multilingual, encoded or low-resource payloads), scenario 5 (payload splitting)} & \cb{3}\,invisible characters\,\tcode{T1.S4} or language switching\,\tcode{T3} or encoding\,\tcode{T2} or emoji substitution\,\tcode{T1.S7} or fragmentation\,\tcode{T4} \\
\xwex{common examples 2-4 and 6, scenarios 2-4 and 6-9 (retrieved content, trusted surfaces, images, memory and RAG, MCP channels)} & \outside{delivery channel, delivery modality, lifecycle} \\
\xwex{common example 7 (fine-tuning interface as gradient oracle)} & \outside{search method, model access} \\
\end{xwrows}\par
\begin{xwrows}
\xwex{mitigation text on attackers mimicking a provenance-marking scheme} & \cb{2}\,structural spoofing \\
\end{xwrows}\par
\end{xwcard}

\needspace{7\baselineskip}
\begin{xwcard}{MITRE ATLAS 2026~\cite{mitreatlas}}{organizes by tactic, then technique}
\xwnote{MITRE ATLAS arranges adversary techniques under tactics. Among its techniques for LLMs, retrieval, and agents, which this card covers, construction appears in three named techniques and in the strategy list of LLM Jailbreak.}
\begin{xwrows}
delay execution of LLM instructions (AML.T0094) & \cb{1}\,conditional trigger\,\tcode{T6.S1} \\
false RAG entry injection (AML.T0071) & \cb{2}\,retrieved-source block\,\tcode{T5.S3} (forged document inside a RAG entry), \xwalso{delivery channel} \\
LLM prompt obfuscation (AML.T0068) & \cb{3}\,standard encoding\,\tcode{T2.S1} or classical cipher\,\tcode{T2.S2} (codes apply to encodings only), \xwalso{delivery rendering, delivery modality} \\
LLM jailbreak (AML.T0054) & \outside{goal, model access}, its strategy bullets are listed as examples \\
LLM prompt injection (AML.T0051) and its direct, indirect, and triggered sub-techniques & \outside{delivery channel} \\
retrieval content crafting, RAG poisoning, prompt infiltration via public-facing application, AI agent tool data poisoning, drive-by compromise & \outside{delivery channel} \\
\end{xwrows}\par
\begin{xwrows}
AI agent tool poisoning (AML.T0110) and its three sub-techniques, AI agent context poisoning (AML.T0080) and its memory and thread sub-techniques & \outside{delivery channel, lifecycle} \\
LLM trusted output components manipulation (AML.T0067) and its citations sub-technique & \outside{goal}, shapes responses to deceive the user \\
agentic resource consumption (AML.T0034.002) & \outside{goal} \\
publish hallucinated entities (AML.T0060) & \outside{lifecycle} \\
AI agent clickbait (AML.T0100) & \outside{target, delivery channel}, baits computer-using agents and AI browsers \\
LLM prompt crafting (AML.T0065) & \outside{search method} \\
\end{xwrows}\par
\begin{xwrows}
extract LLM system prompt, LLM data leakage, AI agent tool invocation, exfiltration and data destruction via AI agent tool invocation & \outside{goal} \\
discover LLM system information (AML.T0069) and its three sub-techniques, LLM prompt self-replication (AML.T0061) & \outside{lifecycle} \\
\xwex{instruction override (AML.T0054 strategy bullet)} & \cb{1}\,override\,\tcode{T2} \\
\xwex{roleplay / persona switching (AML.T0054 strategy bullet)} & \cb{5}\,persona\,\tcode{T1} \\
\xwex{fictionalization and hypotheticals (AML.T0054 strategy bullet)} & \cb{5}\,fiction and hypotheticals\,\tcode{T3} \\
\xwex{separate intent from content (AML.T0054 strategy bullet)} & \cb{1}\,sample request\,\tcode{T7.S1} or composition request\,\tcode{T7.S2} (asks for examples, templates, or analysis) \\
\end{xwrows}\par
\begin{xwrows}
\xwex{multi-turn escalation / Crescendo (AML.T0054 strategy bullet)} & \cb{4}\,commitment\,\tcode{T5} (in each turn), \xwalso{delivery turn structure} \\
\xwex{constrained output formats (AML.T0054 strategy bullet)} & \cb{6}\,schema only\,\tcode{T3.S1} \\
\xwex{obfuscation and transformation (AML.T0054 strategy bullet)} & \cb{3}\,standard encoding\,\tcode{T2.S1} or foreign-language payload\,\tcode{T3.S1} or euphemism\,\tcode{T7.S2} or \cb{6}\,altered answer\,\tcode{T3.S3} (encoding, translation, euphemisms, answer in another language) \\
\xwex{create a high priority objective (AML.T0054 strategy bullet)} & \cb{4}\,security-testing purpose\,\tcode{T4.S2} or harm-prevention purpose\,\tcode{T4.S3} or business purpose\,\tcode{T4.S4} or \cb{1}\,task extension\,\tcode{T3.S1} \\
\xwex{affirmation (AML.T0054 strategy bullet)} & \cb{6}\,response priming\,\tcode{T1} \\
\xwex{algorithmic jailbreak generation, weight manipulation (AML.T0054 description)} & \outside{search method, model access} \\
\end{xwrows}\par
\begin{xwrows}
AI agent configuration, AI agent tools (three techniques), discover LLM hallucinations, gather RAG-indexed targets, modify AI agent configuration, discover AI agent configuration and its four sub-techniques, manipulate user LLM chat history, AI agent, autonomous AI agent communication and its two sub-techniques, AI agent environment reconstruction, deploy AI agent & \outside{lifecycle} \\
AI agent tool, poisoned AI agent tool, modify prompt construction logic & \outside{delivery channel, model access, lifecycle} \\
RAG credential harvesting, credentials from AI agent configuration, RAG databases, AI agent tools, AI agent tool credential harvesting, local AI agent, spearphishing via social engineering LLM, LLM response rendering & \outside{goal, delivery rendering} \\
\end{xwrows}\par
\end{xwcard}

\needspace{7\baselineskip}
\begin{xwcard}{Rehberger 2024~\cite{rehberger2024trustno}}{organizes by impact on confidentiality, integrity, availability}
\xwnote{Rehberger sorts real-world exploit scenarios by their impact on confidentiality, integrity and availability. Construction appears only in conditional injection, ASCII smuggling and one example payload.}
\begin{xwrows}
loss of confidentiality, loss of integrity, loss of availability & \outside{goal} \\
leaking the system prompt, Google Gemini and Google Drive (connecting the user to the attacker via Google Meet), infinite loops or long-running inference & \outside{goal} \\
automatic image rendering and image markdown data exfiltration, automatic unfurling of hyperlinks, clickable hyperlinks, invocation of tools (browsing the web), invocation of tools (Chat with Code plugin) & \outside{goal, target}, exfiltration or actions through the application \\
SpAIware (memory persistent attacks), ZombAIs (Claude Computer Use command and control), persistent output refusal & \outside{goal, lifecycle} \\
Google Docs AI (scams and phishing), refusal via (hidden) instructions & \outside{goal, delivery}, instructions hidden in documents or emails \\
conditional prompt injections & \cb{1}\,conditional trigger\,\tcode{T6.S1} (activates only for a chosen user) \\
\end{xwrows}\par
\begin{xwrows}
ASCII smuggling & \cb{3}\,invisible characters\,\tcode{T1.S4} (also hides exfiltrated output), \xwalso{goal} \\
Terminal DiLLMa (ANSI escape codes) & \outside{target, goal}, terminals interpret the model's output \\
\xwex{example payload for ZombAIs (a web page telling the agent to download and run a file)} & \cb{1}\,plain command\,\tcode{T1} \\
\end{xwrows}\par
\end{xwcard}

\needspace{7\baselineskip}
\begin{xwcard}{BIPIA 2025~\cite{yi2023bipia}}{organizes by application task, attack type, and position}
\xwnote{The benchmark labels indirect attacks by application task, attack type, and position. Its injected sentences are plain commands, some asking for an altered answer.}
\begin{xwrows}
task-irrelevant, task automation, business intelligence, conversational agent, research assistance, sentiment analysis, information retrieval, content creation, learning and tutoring, language translation, programming help & \outside{goal}, another task replaces the user's \\
task-relevant, substitution ciphers, base encoding, reverse text, emoji substitution, rare language translation, alphanumeric substitution, homophonic substitution, misspelling intentionally, anagramming, space removal \& grouping & \cb{6}\,altered answer\,\tcode{T3.S3} (the altered answer is the aim), \xwalso{goal} \\
targeted, information dissemination, marketing \& advertising, entertainment, scams \& fraud, misinformation \& propaganda, instruction, social interaction, persuasion, clickbait, malware distribution & \outside{goal} \\
passive, data eavesdropping, traffic analysis, keylogging, screen scraping, introduce system fingerprinting, cookie theft, memory scanning, dumpster diving, environment variable analysis, device and driver enumeration & \outside{goal}, malicious code added to the answer \\
active, blocking internet connection, corrupting an operating system, encrypting documents and demanding ransom, compromising computers, bringing down hosts and servers, sending out spam emails, crippling critical infrastructures, network propagation, exploiting system vulnerabilities, cryptocurrency mining & \outside{goal}, malicious code added to the answer \\
text attacks, code attacks & \outside{goal, target}, split by payload and application task \\
\end{xwrows}\par
\begin{xwrows}
beginning, middle, end & \outside{delivery position} \\
email QA, web QA, table QA, summarization, code QA & \outside{target}, application tasks \\
\xwex{its injected sentences} & \cb{1}\,added command\,\tcode{T1.S1} (added commands, no other device) \\
\end{xwrows}\par
\end{xwcard}

\needspace{7\baselineskip}
\begin{xwcard}{Guardrail SoK 2026~\cite{wang2026guardrails}}{organizes by turns, then technical paradigm}
\xwnote{The SoK's background typology splits jailbreaks by turns and single-turn attacks by technical paradigm. Only its implicit category is defined by construction.}
\begin{xwrows}
manual jailbreaks, optimization-based jailbreaks, generation-based jailbreaks & \outside{search method} \\
implicit jailbreaks & \cb{3}\,encoding\,\tcode{T2} or fragmentation\,\tcode{T4} or allusion\,\tcode{T7.S3} or foreign-language payload\,\tcode{T3.S1} (substitution cipher, decomposition, clues, translation, encoding) \\
multi-turn jailbreaks & \cb{3}\,fragmentation\,\tcode{T4} or \cb{4}\,commitment\,\tcode{T5} (content of each turn), \xwalso{delivery turn structure} \\
\xwex{AIM, the paper's example of a manual jailbreak} & \cb{5}\,persona\,\tcode{T1} \\
\xwex{GCG's adversarial suffixes, the example of optimization-based jailbreaks} & \cb{3}\,non-semantic sequence\,\tcode{T6}, \xwalso{search method} \\
\end{xwrows}\par
\end{xwcard}

\needspace{7\baselineskip}
\begin{xwcard}{B-I-P model 2025~\cite{trustauth2025bip}}{organizes by belief, intention, and permission stages}
\xwnote{The B-I-P model traces agent attacks through belief formation, intent generation, and permission grant. Construction appears only in the forged dialogue templates it cites.}
\begin{xwrows}
cognitive certainty (hallucination and errors), semantic consistency (adversarial perturbation) & \outside{failure mode}, belief-formation dimensions \\
data provenance (context contamination), tool provenance (supply chain poisoning), return value safety (recursive injection) & \outside{delivery channel, trust} \\
data confidentiality, operational irreversibility, cascade impact scope & \outside{goal}, permission-grant risk dimensions \\
user query, system prompt, tool return, agent return, local resource, protocol metadata & \outside{delivery channel}, victim components in its coding scheme \\
\xwex{the forged dialogue templates it cites (ChatInject)} & \cb{2}\,chat-template token\,\tcode{T2.S1} \,+\, fabricated dialogue\,\tcode{T4} \\
\end{xwrows}\par
\end{xwcard}

\needspace{7\baselineskip}
\begin{xwcard}{Yu et al. 2024~\cite{yu2024dontlisten}}{organizes by shared design principle of jailbreak prompts}
\xwnote{It codes jailbreak prompts by the design principle they share, giving five categories and ten patterns that are mostly pretexts, personas, and output forms.}
\begin{xwrows}
disguised intent, research and testing, joking pretext & \cb{4}\,pretext\,\tcode{T4} (research, testing, or humor as purpose) \\
role play & \cb{5}\,persona\,\tcode{T1} or fiction and hypotheticals\,\tcode{T3} \\
defined persona & \cb{5}\,persona\,\tcode{T1} \\
imagined scenario & \cb{5}\,fiction and hypotheticals\,\tcode{T3} \\
structured response & \cb{6}\,forced format\,\tcode{T3} or response priming\,\tcode{T1} or \cb{3}\,surrogate syntax\,\tcode{T8} \\
language translation & \cb{6}\,altered answer\,\tcode{T3.S3} (output in an uncommon language) \\
\end{xwrows}\par
\begin{xwrows}
text continuation & \cb{6}\,response priming\,\tcode{T1} \\
program execution & \cb{3}\,pseudocode\,\tcode{T8.S2} \\
virtual AI simulation & \cb{5}\,persona\,\tcode{T1} \\
superior mode & \cb{5}\,special mode\,\tcode{T1.S2} \\
opposite mode & \cb{5}\,split persona\,\tcode{T1.S4} or unrestricted AI\,\tcode{T1.S1} \\
alternate model & \cb{5}\,unrestricted AI\,\tcode{T1.S1} \\
\end{xwrows}\par
\begin{xwrows}
hybrid strategies & \cb{3}\,obfuscation or \cb{4}\,social engineering or \cb{5}\,scenario framing or \cb{6}\,output coercion (two or more categories above combined) \\
direct query, initial direct input & \cb{1}\,plain command\,\tcode{T1} (user-study codebook) \\
minimal modification & \cb{1}\,override\,\tcode{T2} (its example asks to override restrictions) \\
resemble existing prompts & \cb{4}\,pretext\,\tcode{T4} or \cb{5}\,scenario framing (reuses the Table 1 categories) \\
AI-assisted prompt design, model as co-designer, model for prompt engineering, model as proxy & \outside{search method} \\
\end{xwrows}\par
\end{xwcard}

\needspace{7\baselineskip}
\begin{xwcard}{MHJ 2024~\cite{li2024mhj}}{organizes by red-teaming tactic}
\xwnote{It labels each human multi-turn jailbreak with one dominant tactic out of seven. All tactics but one describe construction, and that one is multi-turn escalation.}
\begin{xwrows}
direct request & \cb{1}\,plain command\,\tcode{T1} \\
echoing & \cb{6}\,fixed text\,\tcode{T3.S2} or \cb{1}\,transformation object\,\tcode{T3.S3} (replicate, confirm, or expand a statement) \\
hidden intention streamline & \cb{4}\,foot-in-the-door\,\tcode{T5.S1} or \cb{3}\,fragmentation\,\tcode{T4} (escalation or split across turns), \xwalso{delivery turn structure} \\
injection & \cb{6}\,response priming\,\tcode{T1} or forced format\,\tcode{T3} or \cb{1}\,plain command\,\tcode{T1} or \cb{4}\,consent claim\,\tcode{T1.S6} \\
dialogue injection, instruction injection & \cb{6}\,required opening\,\tcode{T1.S1} or forced format\,\tcode{T3} (phrases or content the response must include) \\
mandate, command & \cb{1}\,plain command\,\tcode{T1} or importance flag\,\tcode{T4.S2} or \cb{6}\,forced format\,\tcode{T3} \\
\end{xwrows}\par
\begin{xwrows}
permission & \cb{4}\,consent claim\,\tcode{T1.S6} or \cb{1}\,exception clause\,\tcode{T2.S5} \\
obfuscation & \cb{3}\,obfuscation \\
crowding & \cb{3}\,overload\,\tcode{T5} \\
stylized input, encoded/encrypted input, foreign language & \cb{3}\,encoding\,\tcode{T2} or language switching\,\tcode{T3} (foreign language is T3.S1) \\
synonyms & \cb{3}\,synonym swap\,\tcode{T7.S1} or allusion\,\tcode{T7.S3} or \cb{1}\,meaning remap\,\tcode{T6.S5} (a harmless stand-in word) \\
output format, stylized output, literary elements manipulation & \cb{6}\,forced format\,\tcode{T3} (the latter two are style mandates) \\
\end{xwrows}\par
\begin{xwrows}
requested output & \cb{6}\,style mandate\,\tcode{T3.S5} or \cb{1}\,composition request\,\tcode{T7.S2} (an email, blog post, and the like) \\
debate, outside sources & \cb{1}\,proxy request\,\tcode{T7} or \cb{6}\,style mandate\,\tcode{T3.S5} \\
splitting & \cb{5}\,split persona\,\tcode{T1.S4} (good AI and bad AI) \\
subtraction & \cb{6}\,refusal suppression\,\tcode{T2} (remove disclaimers or warnings) \\
request framing & \cb{5}\,scenario framing or \cb{4}\,social engineering \\
framing as code, code input, code output & \cb{3}\,pseudocode\,\tcode{T8.S2} or \cb{1}\,composition request\,\tcode{T7.S2} (input as code, or code requested) \\
\end{xwrows}\par
\begin{xwrows}
fictionalization, recursion, narration, hypothetical, opposite day/mirror world & \cb{5}\,fiction and hypotheticals\,\tcode{T3} (narration S1, hypothetical S2, opposite day S3) \\
allegory & \cb{3}\,allusion\,\tcode{T7.S3} \\
appeal to authority & \cb{5}\,fiction and hypotheticals\,\tcode{T3} \,+\, \cb{4}\,authority claim\,\tcode{T1} (fiction reinforced by authority) \\
false premise & \cb{4}\,false premise\,\tcode{T7.S8} \\
requesting for beneficial purposes & \cb{4}\,pretext\,\tcode{T4} \\
roleplay (persona creation), slippery slope character & \cb{5}\,persona\,\tcode{T1} (slippery slope character is S3) \\
\end{xwrows}\par
\begin{xwrows}
urgency & \cb{4}\,urgency and threat\,\tcode{T3} \\
emotional appeal/manipulation & \cb{4}\,emotional appeal\,\tcode{T2} \\
\end{xwrows}\par
\end{xwcard}

\needspace{7\baselineskip}
\begin{xwcard}{Inie et al. 2025~\cite{inie2025summon}}{organizes by red teamers' strategies and techniques}
\xwnote{It organizes practitioners' red-teaming approaches into five categories, twelve strategies, and 35 techniques. Most techniques are construction, while most stratagems are search tactics.}
\begin{xwrows}
language, code \& encode & \cb{3}\,obfuscation (code \& encode is T8 or T2) \\
Base64, ROT13 & \cb{3}\,encoding\,\tcode{T2} \\
iPython & \cb{5}\,terminal or interpreter\,\tcode{T2.S1} (model told it can run iPython) \\
SQL & \cb{1}\,composition request\,\tcode{T7.S2} or \cb{6}\,schema only\,\tcode{T3.S1} (SQL that fills a table of content) \\
matrices, transformer translatable tokens & \cb{3}\,non-semantic sequence\,\tcode{T6} (glitch tokens, with matrices the closest match) \\
stop sequences & \cb{2}\,boundary token\,\tcode{T2.S2} or section separator\,\tcode{T1.S2} (listed under two strategies) \\
\end{xwrows}\par
\begin{xwrows}
prompt injection, ignore previous instructions & \cb{1}\,override\,\tcode{T2} \\
strong arm attack & (\cb{1}\,supersession\,\tcode{T2.S3} \,+\, capitals and symbols\,\tcode{T4.S1}) or (\cb{4}\,insider identity\,\tcode{T1.S1} \,+\, \cb{1}\,capitals and symbols\,\tcode{T4.S1}) ('ADMIN OVERRIDE' in capitals) \\
stylizing, formal language, synonymous language & \cb{3}\,rewording\,\tcode{T7} (synonyms S1, formal register closest) \\
servile language & \cb{6}\,required opening\,\tcode{T1.S1} or \cb{4}\,emotional appeal\,\tcode{T2} (e.g., 'gladly', unclear from the text) \\
capitalizing & \cb{1}\,capitals and symbols\,\tcode{T4.S1} \\
give examples & \cb{1}\,worked examples\,\tcode{T5.S4} or \cb{6}\,pattern continuation\,\tcode{T1.S3} \\
\end{xwrows}\par
\begin{xwrows}
rhetoric, persuasion \& manipulation & \cb{4}\,social engineering \\
distraction & \cb{3}\,distractors\,\tcode{T5.S5} or \cb{1}\,transformation object\,\tcode{T3.S3} (e.g., a translation task) \\
escalating & \cb{4}\,foot-in-the-door\,\tcode{T5.S1} (in each turn), \xwalso{delivery turn structure} \\
reverse psychology & \cb{4}\,pretext\,\tcode{T4} (prompter framed as doing the right thing) \\
Socratic questioning, identity characteristics, social hierarchies & \outside{goal}, topic choice to expose bias \\
possible worlds & \cb{5}\,simulation\,\tcode{T2} or fiction and hypotheticals\,\tcode{T3} \\
\end{xwrows}\par
\begin{xwrows}
emulations, unreal computing & \cb{5}\,simulation\,\tcode{T2} (unreal computing emulates Linux (S1)) \\
world building, opposite world & \cb{5}\,alternate world\,\tcode{T3.S3} \\
scenarios & \cb{5}\,fiction and hypotheticals\,\tcode{T3} or \cb{4}\,emergency\,\tcode{T3.S2} or \cb{5}\,special mode\,\tcode{T1.S2} (gladiator fight, emergency, training mode) \\
fictionalizing & \cb{5}\,scenario framing or \cb{6}\,style mandate\,\tcode{T3.S5} \\
switching genres, poetry, forum posts & \cb{6}\,style mandate\,\tcode{T3.S5} or \cb{3}\,verse or genre form\,\tcode{T7.S5} or \cb{1}\,composition request\,\tcode{T7.S2} (genre of the requested output) \\
games & \cb{5}\,game or challenge\,\tcode{T4} \\
\end{xwrows}\par
\begin{xwrows}
re-storying, goal hijacking & \cb{1}\,task extension\,\tcode{T3.S1} or \cb{5}\,story or script\,\tcode{T3.S1} (narrative redirected from within) \\
roleplaying, DAN (Do Anything Now), personas & \cb{5}\,persona\,\tcode{T1} (DAN is an unrestricted AI (S1)) \\
claim authority & \cb{4}\,expert credentials\,\tcode{T1.S2} \\
stratagems, scattershot, regenerate response, clean slate, meta-prompting, ask for examples & \outside{search method}, resampling or attack ideas from the model \\
changing temperature & \outside{model access} \\
perspective-shifting & \cb{5}\,hypothetical\,\tcode{T3.S2} or unrestricted AI\,\tcode{T1.S1} ('what if you didn't have this restriction?') \\
\end{xwrows}\par
\end{xwcard}

\needspace{7\baselineskip}
\begin{xwcard}{Lin et al. 2025~\cite{lin2025achilles}}{organizes by model capability exploited}
\xwnote{It groups jailbreak strategies by the model capability they exploit, tied to training stages. Three of its four categories are construction, and one is model access.}
\begin{xwrows}
completion compliance & \cb{6}\,response priming\,\tcode{T1} or \cb{2}\,structural spoofing or \cb{1}\,instruction manipulation (suffixes, separators, demonstrations) \\
affirmative suffixes & \cb{6}\,response priming\,\tcode{T1} (e.g., appending 'Sure, here is') \\
context switching & \cb{2}\,section separator\,\tcode{T1.S2} or \cb{1}\,instruction nullification\,\tcode{T2.S1} or replacement command\,\tcode{T1.S2} (separators, ignore orders, task switches) \\
in-context learning & \cb{2}\,many-shot demonstrations\,\tcode{T4.S2} or \cb{1}\,worked examples\,\tcode{T5.S4} (labeled or unlabeled demonstrations) \\
instruction indirection & \cb{3}\,rewording\,\tcode{T7} or \cb{6}\,output coercion or \cb{5}\,simulation\,\tcode{T2} or fiction and hypotheticals\,\tcode{T3} \\
input euphemisms & \cb{3}\,rewording\,\tcode{T7} or mid-sentence switch\,\tcode{T3.S2} (partial translation is T3.S2) \\
\end{xwrows}\par
\begin{xwrows}
output constraints & \cb{6}\,forced format\,\tcode{T3} or refusal suppression\,\tcode{T2} \\
virtual simulation & \cb{5}\,simulation\,\tcode{T2} or fiction and hypotheticals\,\tcode{T3} or (\cb{3}\,pseudocode\,\tcode{T8.S2} \,+\, variable concatenation\,\tcode{T4.S1}) (scenario or program-execution simulation) \\
generalization glide & \cb{3}\,language switching\,\tcode{T3} or encoding\,\tcode{T2} or \cb{5}\,persona\,\tcode{T1} or \cb{4}\,social engineering \\
languages & \cb{3}\,language switching\,\tcode{T3} \\
ciphers & \cb{3}\,encoding\,\tcode{T2} or ASCII art\,\tcode{T1.S5} (ASCII art is T1.S5) \\
personification & \cb{5}\,persona\,\tcode{T1} or \cb{4}\,social engineering (role play or psychological manipulation) \\
\end{xwrows}\par
\begin{xwrows}
model manipulation, decoding manipulation, activations manipulation, model fine-tuning & \outside{model access} \\
\xwex{program execution simulation (keyword split into code variables)} & \cb{3}\,pseudocode\,\tcode{T8.S2} \,+\, variable concatenation\,\tcode{T4.S1} \\
\end{xwrows}\par
\end{xwcard}

\needspace{7\baselineskip}
\begin{xwcard}{Dziemian et al. 2026~\cite{dziemian2026vulnerable}}{organizes by strategies in competition submissions}
\xwnote{It lists 27 strategies, distilled from LLM-tagged competition submissions against agents, that describe how indirect injections are written. Nearly all are construction.}
\begin{xwrows}
fake chain of thought & \cb{2}\,assistant turn\,\tcode{T4.S1} \\
request to disable critical and safety-oriented thoughts & \cb{6}\,reasoning ban\,\tcode{T2.S6} \\
offer reward and punishment & \cb{4}\,incentive offer\,\tcode{T7.S7} or threat\,\tcode{T3.S4} \\
fake syntax and delimiters, fake LLM system prompt & \cb{2}\,delimiters\,\tcode{T1} or control tokens\,\tcode{T2} or markup element\,\tcode{T3.S2} (fake boundaries, system labels, or tokens) \\
get agent to follow scripted message & \cb{6}\,fixed text\,\tcode{T3.S2} \\
user provided malicious tool call string & \cb{2}\,assistant turn\,\tcode{T4.S1} or \cb{3}\,pseudocode\,\tcode{T8.S2} (tool-call syntax the agent may run) \\
\end{xwrows}\par
\begin{xwrows}
frame as simulation, game, or fictional context & \cb{5}\,persona\,\tcode{T1} or consequence-free sandbox\,\tcode{T2.S3} or fiction and hypotheticals\,\tcode{T3} or game or challenge\,\tcode{T4} \\
fake user and assistant messages & \cb{2}\,fabricated dialogue\,\tcode{T4} \\
chain harmful shell command with benign commands & \cb{3}\,pseudocode\,\tcode{T8.S2} \,+\, \cb{1}\,routine step\,\tcode{T3.S2} (malicious command chained among benign ones) \\
forge tool or service outputs, fake error state and recovery procedure & \cb{2}\,counterfeit tool results\,\tcode{T5} (error states are S2) \\
excessive repetition & \cb{1}\,repetition\,\tcode{T4.S3} \\
argue valid justification & \cb{4}\,pretext\,\tcode{T4} or argument and obligation\,\tcode{T7} \\
\end{xwrows}\par
\begin{xwrows}
fake social proof & \cb{4}\,social proof\,\tcode{T6} \\
encode or obfuscate text & \cb{3}\,encoding\,\tcode{T2} or invisible characters\,\tcode{T1.S4} (CSS hiding is rendering-level delivery), \xwalso{delivery rendering} \\
explicitly override system rules & \cb{1}\,override\,\tcode{T2} \\
embed in legitimate context & \cb{1}\,task embedding\,\tcode{T3} \\
impersonate authority or user & \cb{4}\,insider identity\,\tcode{T1.S1} or requester impersonation\,\tcode{T1.S4} \\
overwrite agent persona & \cb{5}\,persona\,\tcode{T1} \\
\end{xwrows}\par
\begin{xwrows}
claim urgency or emergency & \cb{4}\,urgency and threat\,\tcode{T3} \\
conditional logic and setting triggers & \cb{1}\,conditional trigger\,\tcode{T6.S1} \\
assert behavior is necessary for downstream benign goal & \cb{4}\,false premise\,\tcode{T7.S8} or \cb{1}\,routine step\,\tcode{T3.S2} (a false prerequisite for the task) \\
command silent or hidden execution & \cb{6}\,forced format\,\tcode{T3} (closest match, as the response must hide the action) \\
redefine real words & \cb{1}\,meaning remap\,\tcode{T6.S5} \\
multilingual text & \cb{3}\,language switching\,\tcode{T3} \\
\end{xwrows}\par
\begin{xwrows}
emotional manipulation & \cb{4}\,emotional appeal\,\tcode{T2} \\
\xwex{its five universal attack clusters (holodeck, protocol override, visual ruleset injection, alignment matrix, context hijacking)} & (\cb{5}\,simulation\,\tcode{T2} \,+\, \cb{2}\,structured-data disguise\,\tcode{T3}) or (\cb{4}\,document guise\,\tcode{T1.S5} \,+\, \cb{1}\,policy revocation\,\tcode{T2.S4}) or (emphasis\,\tcode{T4} \,+\, supersession\,\tcode{T2.S3}) or \cb{5}\,special mode\,\tcode{T1.S2} or \cb{1}\,added command\,\tcode{T1.S1} \\
\end{xwrows}\par
\end{xwcard}

\needspace{7\baselineskip}
\begin{xwcard}{Khodayari et al. 2026~\cite{khodayari2026wild}}{organizes by objective, target, placement, and technique}
\xwnote{It codes web-embedded injections by objective, targeted agent, and placement, and tags eight prompting techniques, which are the only construction part.}
\begin{xwrows}
instruction override & \cb{1}\,override\,\tcode{T2} \\
jailbreak framing & \cb{1}\,policy revocation\,\tcode{T2.S4} or \cb{5}\,special mode\,\tcode{T1.S2} (removes restrictions or enables unrestricted modes) \\
output constraint, content injection & \cb{6}\,forced format\,\tcode{T3} \\
authority pressure & \cb{4}\,authority claim\,\tcode{T1} or urgency and threat\,\tcode{T3} \\
conditional targeting & \cb{1}\,direct aside\,\tcode{T5.S2} or conditional trigger\,\tcode{T6.S1} ('if you are an AI' phrasing) \\
role playing, identity rewrite & \cb{5}\,persona\,\tcode{T1} \\
\end{xwrows}\par
\begin{xwrows}
offensive, defensive, underspecified, system disruption or degradation, reputation manipulation, data exfiltration, data protection, AI bot identification, generic content override & \outside{goal} \\
garbage injection (corruption), model self-modification, command injection, constraint removal, content / product promotion, source citation forcing, positive review forcing, SEO backlink injection, job candidate promotion, reveal system prompt, leak secrets, personal information, copyright, social media bots, challenge-response injection, honeypot & \outside{goal}, objective subcategories \\
data scraping agent, search engine agent, customer support agent, HR screening agent, task executor agent & \outside{target} \\
page response headers, page response body, site-level resources & \outside{delivery channel}, injection surfaces \\
custom header fields, structured data, page metadata, comments, HTML elements and attributes & \outside{delivery position}, embedding mechanisms \\
rendering suppression, visual obfuscation, spatial hiding, layer-based hiding, display \& visibility, small-sized elements, clipping \& masking, text size manipulation, color \& contrast, opacity manipulation, off-screen positioning, viewport visibility, occlusion, z-index manipulation, non-visible by nature, non-visible, visible & \outside{delivery rendering}, hiding techniques and visibility classes \\
\end{xwrows}\par
\begin{xwrows}
\xwex{AI-addressed fields such as X-AI headers and <meta name=``ai-directive''>} & \cb{2}\,metadata field\,\tcode{T3.S4} \\
\end{xwrows}\par
\end{xwcard}

\needspace{7\baselineskip}
\begin{xwcard}{Palo Alto Networks 2025~\cite{paloalto2025securing}}{organizes by technique crossed with impact}
\xwnote{Palo Alto Networks crosses four technique categories with four impact categories. Its technique side names obfuscation and social engineering but mixes in search methods and goals.}
\begin{xwrows}
goal hijacking, guardrail bypass, information leakage, infrastructure attack & \outside{goal} \\
prompt engineering & \cb{1}\,instruction manipulation or \cb{2}\,structural spoofing or \cb{3}\,obfuscation or \cb{5}\,scenario framing or \cb{6}\,output coercion (by member approach) \\
objective manipulation & \cb{1}\,override\,\tcode{T2} or plain command\,\tcode{T1} \\
repeated token, glitch tokens & \cb{3}\,non-semantic sequence\,\tcode{T6} \\
prompt leakage, leak replay, remote code execution, exposure of instructions and tool schemas & \outside{goal} \\
payload splitting & \cb{3}\,fragmentation\,\tcode{T4} \\
\end{xwrows}\par
\begin{xwrows}
system mode & \cb{5}\,special mode\,\tcode{T1.S2} or terminal or interpreter\,\tcode{T2.S1} \\
ReNeLLM & \cb{3}\,rewording\,\tcode{T7} \,+\, \cb{1}\,task embedding\,\tcode{T3} (rewriting plus nesting in a task) \\
fuzzing, AutoDAN, prompt automatic iterative response (PAIR) & \outside{search method} \\
adversarial suffix/prefix & \cb{1}\,added command\,\tcode{T1.S1} (instruction placed before or after input), \xwalso{delivery position} \\
few-shots (many-shots) & \cb{1}\,worked examples\,\tcode{T5.S4} or \cb{2}\,many-shot demonstrations\,\tcode{T4.S2} \\
repeat-instruction attacks & \cb{6}\,fixed text\,\tcode{T3.S2} or length mandate\,\tcode{T3.S6}, \xwalso{goal} \\
\end{xwrows}\par
\begin{xwrows}
social engineering & \cb{4}\,social engineering or \cb{5}\,scenario framing (by member approach, also holding F6) \\
crescendo & \cb{4}\,commitment\,\tcode{T5} (in each turn), \xwalso{delivery turn structure} \\
deceptive delight & \cb{5}\,story or script\,\tcode{T3.S1} (unsafe topic hidden in a story), \xwalso{delivery turn structure} \\
bad likert judge & \cb{5}\,job persona\,\tcode{T1.S6} \,+\, \cb{1}\,sample request\,\tcode{T7.S1}, \xwalso{delivery turn structure} \\
indirect reference & \cb{1}\,proxy request\,\tcode{T7} or \cb{3}\,allusion\,\tcode{T7.S3} \\
skeleton key & \cb{1}\,exception clause\,\tcode{T2.S5} or policy revocation\,\tcode{T2.S4} (adds a warn-instead-of-refuse rule) \\
\end{xwrows}\par
\begin{xwrows}
storytelling, character roleplay, hypothetical scenario & \cb{5}\,story or script\,\tcode{T3.S1} or persona\,\tcode{T1} or hypothetical\,\tcode{T3.S2} (by member approach) \\
persuasion & \cb{4}\,argument and obligation\,\tcode{T7} \\
grandma appeals & \cb{5}\,persona\,\tcode{T1} \,+\, \cb{4}\,sympathy\,\tcode{T2.S1} \\
refusal suppression, affirmative response prefix & \cb{6}\,refusal suppression\,\tcode{T2} or required opening\,\tcode{T1.S1} \\
obfuscation, encoding schemes, cipher text, flip text, prompts in low-resourced languages, special characters, ASCII text & \cb{3}\,standard encoding\,\tcode{T2.S1} or classical cipher\,\tcode{T2.S2} or transposition\,\tcode{T2.S5} or foreign-language payload\,\tcode{T3.S1} or emoji substitution\,\tcode{T1.S7} or invisible characters\,\tcode{T1.S4} or ASCII art\,\tcode{T1.S5} (by member approach) \\
output obfuscation & \cb{6}\,altered answer\,\tcode{T3.S3} (named in Table 2, not defined) \\
\end{xwrows}\par
\begin{xwrows}
knowledge poisoning, backdoor & \outside{delivery channel, lifecycle, model access}, category and approach share a name \\
memory corruption & \outside{lifecycle, delivery channel} \\
direct function exploitation & \outside{target, goal} \\
\end{xwrows}\par
\end{xwcard}

\needspace{7\baselineskip}
\begin{xwcard}{Unit 42 2026~\cite{unit42_2026fooling}}{organizes by attacker intent and payload engineering}
\xwnote{Unit 42 classifies web-based indirect injections by attacker intent and by payload engineering, where construction appears in its jailbreak methods beside delivery-level concealment.}
\begin{xwrows}
attacker intent, low severity, medium severity, high severity, critical severity (with their 13 intents) & \outside{goal} \\
payload engineering & \cb{3}\,obfuscation or \cb{2}\,structural spoofing or \cb{4}\,social engineering (delivery methods and jailbreak methods), \xwalso{delivery rendering} \\
prompt delivery methods, visual concealment, zero-sizing, off-screen positioning, CSS rendering suppression, transparency, camouflage, obfuscation (the delivery method), XML/SVG encapsulation, HTML attribute cloaking, dynamic execution, visible plaintext & \outside{delivery rendering} \\
URL string manipulation & \outside{delivery channel} \\
jailbreak methods & \cb{3}\,obfuscation or \cb{2}\,structural spoofing or \cb{4}\,social engineering (by child method) \\
instruction obfuscation, invisible characters, homoglyph substitution, payload splitting, garbled text, encoding, HTML entity encoding, Base-N encoding, URL encoding, nested encoding & \cb{3}\,invisible characters\,\tcode{T1.S4} or homoglyphs\,\tcode{T1.S3} or fragmentation\,\tcode{T4} or case and typo noise\,\tcode{T1.S6} or standard encoding\,\tcode{T2.S1} (by child method) \\
\end{xwrows}\par
\begin{xwrows}
Unicode bi-directional override & \cb{3}\,invisible characters\,\tcode{T1.S4} (reverses the displayed text), \xwalso{delivery rendering} \\
semantic tricks & \cb{3}\,language switching\,\tcode{T3} or \cb{2}\,structured-data disguise\,\tcode{T3} or \cb{4}\,social engineering (by child trick) \\
multilingual instructions & \cb{3}\,foreign-language payload\,\tcode{T3.S1} \\
JSON/syntax injection & \cb{2}\,structured-data disguise\,\tcode{T3} (closes the JSON, forges fields) \\
social engineering & \cb{4}\,social engineering \\
\xwex{authority cues (god mode, developer mode) and DAN role-play named in its social-engineering description} & \cb{5}\,special mode\,\tcode{T1.S2} or unrestricted AI\,\tcode{T1.S1} \\
\end{xwrows}\par
\end{xwcard}

\needspace{7\baselineskip}
\begin{xwcard}{Microsoft Prompt Shields 2024~\cite{microsoft2024promptshields}}{organizes by attack channel, then attack subtype}
\xwnote{Prompt Shields splits attacks by channel and lists four construction-level subtypes for both channels, plus six goal subtypes for document attacks.}
\begin{xwrows}
user prompt attacks, document attacks & \outside{delivery channel} \\
attempt to change system rules & \cb{1}\,override\,\tcode{T2} or \cb{5}\,unrestricted AI\,\tcode{T1.S1} (listed for both channels) \\
embedding a conversation mockup to confuse the model & \cb{2}\,fabricated dialogue\,\tcode{T4} (listed for both channels) \\
role-play & \cb{5}\,persona\,\tcode{T1} (listed for both channels) \\
encoding attacks & \cb{3}\,character perturbation\,\tcode{T1} or encoding\,\tcode{T2} or rewording\,\tcode{T7} or language switching\,\tcode{T3} (character, cipher, style or language variants) \\
manipulated content, intrusion, information gathering, availability, fraud, malware & \outside{goal} \\
\end{xwrows}\par
\end{xwcard}

\needspace{7\baselineskip}
\begin{xwcard}{Arcanum 2026~\cite{haddix2025arcanum}}{organizes by intent, technique, evasion and input pillars}
\xwnote{Arcanum files each node under one of four pillars, intents, techniques, evasions and inputs, and construction appears across its technique and evasion pillars.}
\begin{xwrows}
intents (the pillar and all 27 intents) & \outside{goal} \\
inputs (the pillar and all 12 inputs) & \outside{delivery channel, delivery modality} \\
techniques & \cb{1}\,instruction manipulation or \cb{2}\,structural spoofing or \cb{3}\,obfuscation or \cb{4}\,social engineering or \cb{5}\,scenario framing or \cb{6}\,output coercion (pillar, see member rows) \\
direct request (plain prompting), contradiction, reorientation, rule addition, special-case exception, reiteration, conditional / trigger-gated payload (sleeper), inversion & \cb{1}\,plain command\,\tcode{T1} or override\,\tcode{T2} or execution rules\,\tcode{T6} or repetition\,\tcode{T4.S3} or conditional trigger\,\tcode{T6.S1} or inverted request\,\tcode{T7.S3} (by member technique) \\
end sequences, many-shot jailbreaking, history fabrication (fake assistant turn), policy-file framing (Policy Puppetry), special-token injection, fake completion, chain-of-thought spoofing, tool-call spoofing & \cb{2}\,section separator\,\tcode{T1.S2} or control tokens\,\tcode{T2} or runtime config\,\tcode{T3.S1} or fabricated dialogue\,\tcode{T4} or counterfeit tool results\,\tcode{T5} (by member technique) \\
binary streams, cognitive overload, figurative language, shortcuts, variable expansion, glitch / anomalous tokens, context overflow (window flooding), reasoning dilution (CoT hijacking), masked-word reconstruction (SATA), induced hallucination & \cb{3}\,standard encoding\,\tcode{T2.S1} or overload\,\tcode{T5} or rewording\,\tcode{T7} or variable concatenation\,\tcode{T4.S1} or pseudocode\,\tcode{T8.S2} or non-semantic sequence\,\tcode{T6} or fragmentation\,\tcode{T4} (by member technique) \\
\end{xwrows}\par
\begin{xwrows}
anti-harm coercion, urgency, persuasion (social-engineering levers), authority impersonation, self-persuasion (self-generated rationalization), fake-citation grounding (DarkCite), agentic compliance momentum (foot-in-the-door), memory exploitation & \cb{4}\,pretext\,\tcode{T4} or argument and obligation\,\tcode{T7} or urgency and threat\,\tcode{T3} or \cb{4}\,social engineering or authority claim\,\tcode{T1} or commitment\,\tcode{T5} or cited evidence\,\tcode{T7.S10} or consent claim\,\tcode{T1.S6} (by member technique) \\
act as interpreter, narrative injection (aka framing), competition, tense reformulation (past / future tense), evaluator-role abuse (Bad Likert Judge), distraction sandwich (Deceptive Delight) & \cb{5}\,terminal or interpreter\,\tcode{T2.S1} or fiction and hypotheticals\,\tcode{T3} or persona\,\tcode{T1} or game or challenge\,\tcode{T4} or era setting\,\tcode{T3.S5} or (job persona\,\tcode{T1.S6} \,+\, \cb{1}\,sample request\,\tcode{T7.S1}) or (\cb{5}\,story or script\,\tcode{T3.S1} \,+\, \cb{3}\,distractors\,\tcode{T5.S5}) (by member technique) \\
anti-refusal, chunking, priming, truncated instructions, output priming (prefix injection), thinking-mode manipulation (reasoning-budget steering), structured-output coercion (constrained decoding) & \cb{6}\,refusal suppression\,\tcode{T2} or partial answer\,\tcode{T3.S4} or required opening\,\tcode{T1.S1} or forced format\,\tcode{T3} or reasoning ban\,\tcode{T2.S6} or thought seeding\,\tcode{T1.S4} or schema only\,\tcode{T3.S1} (by member technique) \\
chain of thought introspection, secret probing (oracle extraction) & \cb{1}\,proxy request\,\tcode{T7} or \cb{6}\,thought seeding\,\tcode{T1.S4} (aimed at hidden secrets), \xwalso{goal} \\
link injection & \cb{1}\,fetch and follow\,\tcode{T6.S4} or \cb{3}\,standard encoding\,\tcode{T2.S1} (fetch-and-follow or URL-encoded text) \\
meta prompting & \cb{1}\,proxy request\,\tcode{T7} (asks the model to write the attack), \xwalso{search method} \\
\end{xwrows}\par
\begin{xwrows}
puzzling & \cb{3}\,encoding\,\tcode{T2} or \cb{6}\,pattern continuation\,\tcode{T1.S3} or \cb{5}\,game objective\,\tcode{T4.S1} (cipher, pattern completion or riddle) \\
crescendo (gradual escalation), echo chamber (context poisoning), multi-turn decomposition (sub-query splitting) & \cb{4}\,commitment\,\tcode{T5} or \cb{3}\,fragmentation\,\tcode{T4} (per turn, commitment or fragments), \xwalso{delivery turn structure} \\
gradient-based attacks & \cb{3}\,non-semantic sequence\,\tcode{T6} (optimized token string), \xwalso{search method, model access} \\
best-of-N (augmentation sampling) & \cb{3}\,character perturbation\,\tcode{T1} (sampled character noise), \xwalso{search method} \\
fuzzing-based jailbreak, autonomous strategy discovery (AutoDAN-Turbo) & \outside{search method} \\
abliteration / weight ablation & \outside{model access} \\
\end{xwrows}\par
\begin{xwrows}
tool-definition injection (MCP tool poisoning), retrieval ranking manipulation (RAG poisoning), agent instruction-file injection (rules-file backdoor) & \outside{delivery channel} \\
tool rug pull (TOCTOU mutation), prompt worm (self-replication), Russian doll & \outside{lifecycle, delivery channel} \\
confused deputy (agent authority confusion), tool-preference manipulation (tool squatting), function-call parameter smuggling & \outside{target, trust, delivery channel} \\
spatial byte arrays, steganography, waveforms / frequencies & \outside{delivery modality}, one technique and two evasions \\
evasions & \cb{3}\,obfuscation (pillar, see member rows) \\
A1Z26 number substitution, ancient scripts, Baconian cipher, Base64, binary, bijection learning, Brainfuck / esoteric languages, Braille, cipher, code switching / randomizer, hexadecimal, HTML entities, Japanese scripts, Morse code, NATO phonetic alphabet, phonetic substitution, rail fence cipher, reverse, semaphore, tap code, URL encoding, alternative base encodings (Base32 / Base58 / Base85), numeric code-point encoding (octal / decimal), layered encoding (encoding chains), compression encoding (gzip / zlib), legacy charset confusion (GB18030 / UTF-7) & \cb{3}\,standard encoding\,\tcode{T2.S1} or classical cipher\,\tcode{T2.S2} or phonetic respelling\,\tcode{T2.S4} or transposition\,\tcode{T2.S5} or encoding\,\tcode{T2} (by member evasion) \\
\end{xwrows}\par
\begin{xwrows}
ASCII, bubble / enclosed text, case changing, emoji, fullwidth characters, homoglyphs, invisible text, mathematical Unicode, metacharacter confusion, regional indicators, truncation \& misspelling, spaces / whitespace, splats, small caps / subscript / superscript, strikethrough / underline / overlay, upside down text, vertical text, Wingdings / symbol fonts, Zalgo text & \cb{3}\,ASCII art\,\tcode{T1.S5} or case and typo noise\,\tcode{T1.S6} or emoji substitution\,\tcode{T1.S7} or homoglyphs\,\tcode{T1.S3} or invisible characters\,\tcode{T1.S4} or spacing changes\,\tcode{T1.S2} (by member evasion) \\
acrostics, graph nodes, alternative language, fictional \& constructed languages, synonyms / word substitution, adversarial poetry, symbolic math encoding (MathPrompt), query-language encoding (QueryAttack), code-structure encoding (CodeAttack) & \cb{3}\,positional extraction\,\tcode{T4.S2} or language switching\,\tcode{T3} or classical cipher\,\tcode{T2.S2} or synonym swap\,\tcode{T7.S1} or verse or genre form\,\tcode{T7.S5} or surrogate syntax\,\tcode{T8} or pseudocode\,\tcode{T8.S2} (by member evasion) \\
JSON, XML & \cb{3}\,surrogate syntax\,\tcode{T8} or \cb{2}\,structured-data disguise\,\tcode{T3} (F2 when posing as system data) \\
link smuggling, markdown, ANSI escape concealment (terminal codes) & \outside{delivery rendering} \\
bidirectional text override (Trojan Source) & \cb{3}\,invisible characters\,\tcode{T1.S4} (display order differs from logical order), \xwalso{delivery rendering} \\
delimiter & \cb{6}\,altered answer\,\tcode{T3.S3} or schema only\,\tcode{T3.S1} (formats the answer, not the input) \\
\end{xwrows}\par
\end{xwcard}

\needspace{7\baselineskip}
\begin{xwcard}{Cisco 2025~\cite{cisco2025framework}}{organizes by objective, then technique and sub-technique}
\xwnote{Cisco nests techniques and sub-techniques under attacker objectives. Within prompt injection and jailbreak, construction appears only as a few named sub-techniques.}
\begin{xwrows}
goal hijacking, jailbreak & \outside{goal}, objectives OB-001, OB-002 and technique AITech-2.1 \\
direct prompt injection, indirect prompt injection & \outside{delivery channel}, split by channel \\
instruction manipulation (direct prompt injection), instruction manipulation (indirect prompt injection) & \cb{1}\,override\,\tcode{T2} or replacement command\,\tcode{T1.S2} (override, replace or modify instructions) \\
obfuscation (direct prompt injection), obfuscation (indirect prompt injection), obfuscation (jailbreak) & \cb{3}\,obfuscation (by name, as the report gives no definition) \\
multi-agent prompt injection, multi-agent (indirect prompt injection), multi-agent jailbreak collaboration & \outside{delivery channel} \\
goal manipulation, goal manipulation (models, agents), goal manipulation (tools, prompts, resources) & \outside{goal, delivery channel} \\
\end{xwrows}\par
\begin{xwrows}
multi-modal injection and manipulation, image-text injection, image manipulation, audio command injection, video overlay manipulation & \outside{delivery modality} \\
context manipulation (jailbreak) & \cb{2}\,structural spoofing or \cb{5}\,scenario framing (by name, as the report gives no definition) \\
semantic manipulation (jailbreak) & \cb{3}\,rewording\,\tcode{T7} or \cb{4}\,social engineering (by name, as the report gives no definition) \\
token exploitation (jailbreak) & \cb{2}\,control tokens\,\tcode{T2} or \cb{3}\,non-semantic sequence\,\tcode{T6} (by name, as the report gives no definition) \\
\end{xwrows}\par
\end{xwcard}

\needspace{7\baselineskip}
\begin{xwcard}{SaTML CTF 2024~\cite{debenedetti2024satml}}{organizes by strategies of the winning teams}
\xwnote{The SaTML CTF report describes the strategies that winning teams used to extract a secret, so construction appears only in narrative descriptions and example chats.}
\begin{xwrows}
\xwex{letter alteration} & \cb{3}\,leetspeak\,\tcode{T1.S1} \\
\xwex{synonyms and periphrasis for the secret} & \cb{3}\,rewording\,\tcode{T7} \\
\xwex{ignoring pre-extracted decoy secrets} & \cb{1}\,plain command\,\tcode{T1} or instruction nullification\,\tcode{T2.S1} (decoys are found by probing), \xwalso{search method} \\
\xwex{per-character disclosure with case tags, a Python list of the characters} & \cb{6}\,altered answer\,\tcode{T3.S3} or schema only\,\tcode{T3.S1} \\
\xwex{keeping the answer brief} & \cb{6}\,style mandate\,\tcode{T3.S5} \\
\xwex{benign questions before the request, multi-turn conversations that build context} & \cb{4}\,foot-in-the-door\,\tcode{T5.S1} (often spread over several turns), \xwalso{delivery turn structure} \\
\end{xwrows}\par
\begin{xwrows}
\xwex{opening with ``Sure, here is''} & \cb{6}\,required opening\,\tcode{T1.S1} \\
\xwex{a reply template that begins with a refusal} & \cb{6}\,staged denial\,\tcode{T1.S5} \\
\xwex{an in-context encoding of each character (a word repeated three to five times)} & \cb{6}\,altered answer\,\tcode{T3.S3} \,+\, \cb{1}\,worked examples\,\tcode{T5.S4} \\
\xwex{asking for one or two characters at a time} & \cb{6}\,partial answer\,\tcode{T3.S4} \\
\xwex{replacing patterns that defenses block} & \cb{6}\,altered answer\,\tcode{T3.S3} \\
\xwex{framing the request as help to prevent leakage} & \cb{4}\,harm-prevention purpose\,\tcode{T4.S3} \,+\, \cb{1}\,composition request\,\tcode{T7.S2} (asks for a script containing the secret) \\
\end{xwrows}\par
\begin{xwrows}
\xwex{the buffer fill attack} & \cb{1}\,composition request\,\tcode{T7.S2} (fixing login code yields the secret) \\
\xwex{impersonating a character such as a professor} & (\cb{4}\,expert credentials\,\tcode{T1.S2} \,+\, \cb{1}\,importance flag\,\tcode{T4.S2}) or (\cb{5}\,job persona\,\tcode{T1.S6} \,+\, \cb{1}\,importance flag\,\tcode{T4.S2}) (the source leaves the persona's owner unclear) \\
\xwex{probing the filters as an oracle, checking guesses against the defense} & \outside{search method} \\
\xwex{a successful chat that asks what a program would output} & \cb{3}\,pseudocode\,\tcode{T8.S2} \,+\, \cb{6}\,altered answer\,\tcode{T3.S3} \,+\, schema only\,\tcode{T3.S1} (the program encodes the secret) \\
\xwex{an unsuccessful chat that describes a four-component hypothetical system} & \cb{5}\,hypothetical\,\tcode{T3.S2} \,+\, simulation\,\tcode{T2} \,+\, \cb{4}\,false premise\,\tcode{T7.S8} \,+\, \cb{6}\,altered answer\,\tcode{T3.S3} \,+\, required opening\,\tcode{T1.S1} (ASCII codes are requested in the output) \\
\xwex{attacks a defending team guarded against (ASCII art, character modifications, role plays, games, translations, odd formats, overriding instructions)} & \cb{3}\,ASCII art\,\tcode{T1.S5} or character perturbation\,\tcode{T1} or \cb{5}\,persona\,\tcode{T1} or game or challenge\,\tcode{T4} or \cb{6}\,altered answer\,\tcode{T3.S3} or forced format\,\tcode{T3} or \cb{1}\,override\,\tcode{T2} \\
\end{xwrows}\par
\end{xwcard}

\needspace{7\baselineskip}
\begin{xwcard}{Native label schemes of the ingested datasets}{organizes by various}
\xwnote{Each dataset of Table~\ref{tab:datasets} that carries native attack labels uses its own vocabulary, except that CyberSecEval 2 and 3 share one. All native labels are kept in the \textit{pi\_technique} field.}
\begin{xwrows}
InjecAgent's direct harm and data stealing & \outside{goal} \\
InjecAgent's enhanced setting & \cb{1}\,instruction nullification\,\tcode{T2.S1} \,+\, capitals and symbols\,\tcode{T4.S1} \,+\, importance flag\,\tcode{T4.S2} \\
AgentDojo's and AgentDyn's important message & \cb{1}\,task extension\,\tcode{T3.S1} \,+\, \cb{2}\,markup element\,\tcode{T3.S2} \,+\, \cb{4}\,requester impersonation\,\tcode{T1.S4} \\
AgentDojo's TODO and ignore previous instructions & \cb{1}\,plain command\,\tcode{T1} or instruction nullification\,\tcode{T2.S1} \\
AgentDyn's injection tasks & \outside{goal} \\
ASB's direct, indirect, memory, and plan-of-thought injection and mixed attacks & \outside{delivery, lifecycle} \\
\end{xwrows}\par
\begin{xwrows}
ASB's naive, escape characters, context ignoring, fake completion, and combined, and PIArena's combined & see the Open-Prompt-Injection card \\
SEP's probe position & \outside{delivery position} \\
BrowseSafe's important message, todo, and ignore previous & (\cb{1}\,task extension\,\tcode{T3.S1} \,+\, \cb{2}\,markup element\,\tcode{T3.S2} \,+\, \cb{4}\,requester impersonation\,\tcode{T1.S4}) or (\cb{1}\,plain command\,\tcode{T1} \,+\, importance flag\,\tcode{T4.S2}) or instruction nullification\,\tcode{T2.S1} \\
BrowseSafe's InjecAgent & \cb{1}\,plain command\,\tcode{T1} \,+\, \cb{2}\,metadata field\,\tcode{T3.S4} \\
BrowseSafe's role manipulation, delimiter injection, and social engineering & \cb{5}\,job persona\,\tcode{T1.S6} or \cb{2}\,section separator\,\tcode{T1.S2} or (\cb{4}\,insider identity\,\tcode{T1.S1} \,+\, urgency and threat\,\tcode{T3}) \\
BrowseSafe's indirect hypothetical and multilanguage & \cb{5}\,hypothetical\,\tcode{T3.S2} or \cb{3}\,language switching\,\tcode{T3} \\
\end{xwrows}\par
\begin{xwrows}
BrowseSafe's URL segment and system prompt exfiltration & \outside{delivery channel, goal} \\
BrowseSafe's explicit, indirect, and stealth styles & \cb{1}\,plain command\,\tcode{T1} or override\,\tcode{T2} or routine step\,\tcode{T3.S2} or (task embedding\,\tcode{T3} \,+\, \cb{4}\,official notice\,\tcode{T1.S3}) \\
BrowseSafe's nine injection strategies (hidden elements and rewritten page sections) & \outside{delivery rendering, delivery position} \\
PIArena's phishing, content promotion, access denial, and infrastructure failure & \outside{goal} \\
PIArena's direct attack & \cb{1}\,plain command\,\tcode{T1} \\
PIArena's strategy-based attack & \outside{search method} \\
\end{xwrows}\par
\begin{xwrows}
Open-Prompt-Injection, LLMail-Inject, Gandalf, SaTML CTF, CyberSecEval, PromptInject & see their cards above \\
\end{xwrows}\par
\end{xwcard}
